\chapter{Physical Origins of Holomorphic-Topological Field Theory}
\label{chap:ht-physical-origins}

%================================================================
% SECTION: PHYSICAL ORIGINS
% (from physical_origins.tex)
%================================================================

\section{Physical origins}
\label{sec:physical-origins-overview}
\label{chap:physical-origins}
\label{rem:programme-origins}

String amplitudes, gauge theory partition functions, and conformal
blocks each independently produce bar-cobar structures: the
generating function of $n$-point functions is the bar complex, and
the factorization of amplitudes under degeneration is the coproduct.
This chapter traces the bar-cobar adjunction to its physical sources
and formulates the identifications as precise conjectures.
Theorem-labeled statements are proved; all others are conjectural.

\subsection{4d/2d correspondence algebras}
\label{sec:4d2d}

\subsubsection{The 4d/2d correspondence (BLLPRR)}

\begin{construction}[4d/2d correspondence (BLLPRR)]\label{constr:4d2d}
Starting from a 4d $\mathcal{N} = 2$ superconformal field theory $\mathcal{T}_{4d}$:
\begin{enumerate}[label=(\roman*)]
\item Select the \emph{Schur supercharge} $\mathbb{Q} = Q_1^- + \widetilde{S}^{1\dot{-}}$, a specific nilpotent supercharge in the $\mathcal{N}=2$ superconformal algebra ($\mathbb{Q}^2 = 0$).
\item Pass to the cohomology $H^*(\mathbb{Q})$ of local operators inserted at a fixed plane $\mathbb{R}^2 \subset \mathbb{R}^4$.
\item The resulting algebra $\mathcal{A}[\mathcal{T}_{4d}] = H^*(\mathbb{Q}, \text{local ops})$ carries the structure of a 2d vertex algebra (chiral algebra).
\end{enumerate}

This construction, due to Beem--Lemos--Liendo--Peelaers--Rastelli--van~Rees~\cite{beem-et-al}, produces vertex algebras from 4d SCFTs.
It is distinct from the $\Omega$-background/Nekrasov partition-function
approach used in the AGT correspondence below.
\end{construction}

\begin{theorem}[4d/2d is \texorpdfstring{$\Einf$}{E-infinity}-chiral; \ClaimStatusProvedElsewhere{} \cite{beem-et-al}]\label{thm:4d2d-einf}
The 2d chiral algebra $\mathcal{A}[\mathcal{T}_{4d}]$ obtained from the 4d/2d
correspondence is an $\Einf$-chiral algebra (vertex algebra).

The OPE of the resulting chiral algebra is generically \emph{non-commutative}: for instance, the 4d/2d image of $SU(2)$ $\mathcal{N}=4$ SYM is the small $\mathcal{N}=4$ superconformal vertex algebra at $c = -9$, which contains $\widehat{\mathfrak{sl}}_2$ at level $k = -3/2$ (an admissible, non-integral level; determined by $c = 3k/(k+2)$) as its R-symmetry current subalgebra.  The $\Einf$ structure arises from the full vertex algebra axioms (locality, associativity of OPE), not from commutativity.
\end{theorem}

\begin{example}[Class \texorpdfstring{$\mathcal{S}$}{S}]\label{ex:class-s}
For class $\mathcal{S}$ theories of type $A_{n-1}$ on a Riemann surface $C$:
\[
\mathcal{A}[\mathcal{T}_{A_{n-1}}(C)] = \mathcal{W}^{-n}(\mathfrak{sl}_n)
\]
the W-algebra associated to $\mathfrak{sl}_n$ at the critical level $k = -n = -h^\vee$. At this special level the Sugawara construction degenerates and $\mathcal{W}^{-h^\vee}(\mathfrak{sl}_n)$ coincides with the Feigin--Frenkel center $\mathfrak{z}(\widehat{\mathfrak{g}})$, a commutative vertex algebra (see Feigin--Frenkel \cite{FF} for the center isomorphism and Arakawa \cite{Ara15} for related rationality results). The dependence on the Riemann surface $C$ and its pants decomposition appears at the level of conformal blocks and modules, not the algebra itself.
\end{example}


\subsection{Non-commutative Chern--Simons theory}
\label{sec:nc-chern-simons}
\noindent\emph{All claims in this section are \ClaimStatusConjectured.}

\subsubsection{Chern--Simons as source of chiral algebras}

\begin{construction}[CS boundary algebra]\label{constr:cs-boundary}
3d Chern--Simons theory with gauge group $G$ at level $k$ on $M^3 = \Sigma \times \bR_+$
produces:
\begin{enumerate}[label=(\roman*)]
\item On the boundary $\Sigma$: WZW model at level $k$.
\item The chiral algebra is the affine Kac--Moody algebra $\hat{\mathfrak{g}}_k$.
\end{enumerate}
\end{construction}

\begin{conjecture}[Non-commutative CS; \ClaimStatusConjectured]\label{conj:nc-cs}
Noncommutative Chern--Simons theory on $\bR^3_\theta$ produces
$\Eone$-chiral algebras:
\begin{enumerate}[label=(\roman*)]
\item The boundary theory is a non-commutative WZW model.
\item The OPE is R-twisted with $R$ depending on the non-commutativity parameter $\theta$.
\item At $\theta = 0$, we recover the standard $\Einf$-chiral Kac--Moody.
\end{enumerate}
\end{conjecture}

\begin{remark}[Scope]
The $R$-twisted OPE structure is consistent with Definition~\ref*{def:chiral-ass-operad}, but deriving it from non-commutative CS requires Seiberg--Witten deformation quantization~\cite{SW99}, which lies outside this monograph.
\end{remark}


\subsubsection{\texorpdfstring{Noncommutative geometry and associative / topological $\mathsf{E}_1$ structures}{Noncommutative geometry and associative / topological E1 structures}}

\begin{construction}[NC torus as associative deformation]
\label{constr:nc-torus-e1}
\index{noncommutative torus!topological $E_1$ structure@topological $\mathsf{E}_1$ structure}
The noncommutative torus $T^2_\theta$ ($\theta \in \mathbb{R} \setminus \mathbb{Q}$), with Moyal star product $f \star_\theta g$ deforming $\{x,y\}=\theta$, carries a topological $\Eone$-algebra structure ($VU = e^{2\pi i\theta}UV$, associative but noncommutative).  This is \emph{not} a Beilinson--Drinfeld $\Eone$-chiral algebra on a curve.
\end{construction}

\begin{conjecture}[Koszul duality as Morita equivalence; \ClaimStatusConjectured]
\label{conj:koszul-morita}
\index{Morita equivalence!Koszul duality interpretation}
$\Eone$-Koszul duality for noncommutative tori exchanges the deformation parameter:
\[
\mathcal{A}_\theta^! \;\cong\; \mathcal{A}_{1/\theta}.
\]
This coincides with Rieffel's Morita equivalence $T^2_\theta \sim_{\mathrm{Morita}} T^2_{1/\theta}$
for irrational $\theta$, reinterpreted as $\Eone$-Koszul duality.
The bar complex $\barB(\mathcal{A}_\theta)$ computes the cyclic homology
$\mathrm{HC}_*(\mathcal{A}_\theta)$, which by Connes' theorem equals the de~Rham
cohomology $H^*_{\mathrm{dR}}(T^2)$ --- independent of $\theta$ (topological invariance
of cyclic homology for smooth deformations).
\end{conjecture}

\begin{remark}[Scope]
Rieffel Morita equivalence and Connes' $\mathrm{HC}_*(A_\theta)$ computation are established; the missing step is constructing the topological $\Eone$ bar complex for $\mathcal{A}_\theta$ and verifying its cohomology recovers the Koszul dual at $1/\theta$, which lies outside the algebraic-curve setting.
\end{remark}


\subsection{Gauge theory and D-branes}
\label{sec:gauge-d-branes}
\subsubsection{D-brane vertex algebras}

\begin{construction}[Open string VOA]\label{constr:open-string-voa}
Open strings on D-branes in type IIB produce boundary chiral algebras depending on the brane configuration.
\end{construction}

\begin{conjecture}[D-brane algebras are \texorpdfstring{$\Eone$}{E1}; \ClaimStatusConjectured]\label{conj:dbrane-e1}
Open string vertex algebras on non-trivial D-brane configurations are generically $\Eone$-chiral: Chan--Paton factors give matrix-valued OPE ($\Phi^{ij}(z)\Psi^{k\ell}(w) \sim \delta^{jk}\cdots$), breaking $\Einf$-commutativity.
\end{conjecture}

\begin{remark}[Scope]
Matrix-valued OPE from Chan--Paton factors is compatible with $\Eone$-chiral structure (Definition~\ref*{def:chiral-ass-operad}), but deriving it from open string field theory~\cite{Zwi93} lies outside this monograph.
\end{remark}

\subsubsection{D-brane matrix ordering and the \texorpdfstring{$\Eone$}{E1} Steinberg}
\label{sec:chan-paton-steinberg}

\providecommand{\SteinbergCP}{\mathfrak{S}_b^{\mathrm{CP}}}

\begin{construction}[Chan--Paton Steinberg]\label{constr:chan-paton-steinberg}
For $n$ coincident D-branes, the open string vertex algebra is $\End(\mathbb{C}^n)$-valued.
The Steinberg object of \S\ref{sec:steinberg-from-bar} acquires Chan--Paton dressing:
\[
  \SteinbergCP \;=\; \mathfrak{S}_b \otimes_{\mathbb{C}} \End(\mathbb{C}^n).
\]
The convolution product on $H_*^{\mathrm{BM}}(\SteinbergCP)$ combines two independent structures:
\emph{geometric convolution} from FM collision in the $\mathbb{C}_z$-direction (encoding OPE residues),
and \emph{matrix multiplication} from the Chan--Paton factor $\End(\mathbb{C}^n)$
(encoding open string endpoint labels).
\end{construction}

\begin{proposition}[Matrix ordering forces $\Eone$ structure; \ClaimStatusProvedHere]
\label{prop:chan-paton-e1}
The convolution on $\SteinbergCP$ is $\Eone$ but not $E_\infty$.

The ordered configuration space $\operatorname{Conf}_k^{<}(\mathbb{R})$
in the topological $\mathbb{R}_t$-direction provides a canonical total order on operator
insertions --- this is the $\Eone$-structure of the bar coproduct, not a choice.
For operators $A, B \in \End(\mathbb{C}^n)$ at positions $t_1 < t_2$ in
$\operatorname{Conf}_2^{<}(\mathbb{R})$, the convolution product composes
Chan--Paton factors as $AB$ (left-to-right in the $\mathbb{R}$-ordering).
Since $AB \neq BA$ generically for $n \geq 2$, the product is associative but
non-commutative---i.e.\ $\Eone$ but not $E_\infty$.
The holomorphic $\mathbb{C}_z$-direction is unaffected: OPE coefficients remain
$\Sigma_k$-equivariant on $\operatorname{FM}_k(\mathbb{C})$.
\end{proposition}

\begin{proof}
Associativity of the convolution follows from associativity of both FM composition
(Proposition~\ref*{prop:steinberg-convolution}) and matrix multiplication.
Non-commutativity is immediate: choose $A, B \in \End(\mathbb{C}^n)$ with $[A,B] \neq 0$.
At positions $t_1 < t_2$ the convolution gives $AB$; at $t_2 < t_1$ it gives $BA$.
Since $AB \neq BA$, the $E_\infty$-structure (which requires the transposition
$\sigma \in \Sigma_2$ to act trivially) is broken.  The $\mathbb{C}_z$-factor is
unchanged because FM composition on $\operatorname{FM}_k(\mathbb{C})$ is $\Sigma_k$-equivariant,
and the Chan--Paton tensor factor does not couple to the holomorphic direction.
\end{proof}

\begin{remark}[First-principles origin of $\Eone$-chirality]
\label{rmk:chan-paton-e1-origin}
Proposition~\ref{prop:chan-paton-e1} gives a clean mechanism for why
D-brane algebras are $\Eone$-chiral (Conjecture~\ref{conj:dbrane-e1}):
the Chan--Paton matrix indices create an ordering along the topological $\mathbb{R}_t$-direction
that is invisible in the holomorphic $\mathbb{C}_z$-direction.
For a single brane ($n=1$), $\End(\mathbb{C}^1) = \mathbb{C}$ is commutative and the
$E_\infty$-structure survives---consistent with the fact that abelian open string
vertex algebras are commutative.
\end{remark}


\subsubsection{\texorpdfstring{D-brane categories from $\Eone$-module Koszul duality}{D-brane categories from E1-module Koszul duality}}

\begin{conjecture}[HMS from \texorpdfstring{$\Eone$}{E1}-module Koszul duality; \ClaimStatusConjectured]
\label{conj:hms-koszul}
\index{homological mirror symmetry!Koszul duality interpretation}
\index{D-brane!module Koszul duality}
In string theory, D-branes on a target space $X$ form the derived category
$D^b(\mathrm{Coh}(X))$ (B-branes) or the Fukaya category $\mathrm{Fuk}(X)$
(A-branes).  For a chiral algebra $\mathcal{A}$ on a curve $\Sigma$, the
module category $\mathrm{Mod}(\mathcal{A})$ is the mathematical incarnation
of the category of branes.

Under $\Eone$-Koszul duality, the module bar-cobar functor
$\Phi\colon \mathrm{Mod}(\mathcal{A}) \xrightarrow{\sim}
\mathrm{CoMod}(\mathcal{A}^!)$
of Theorem~\ref*{thm:e1-module-koszul-duality} exchanges module and comodule
categories.  When $X$ is a Calabi--Yau surface and $\mathcal{A}$ is the chiral
algebra of the sigma model on $X$, this functor should recover homological
mirror symmetry (Kontsevich):
\[
D^b\bigl(\mathrm{Coh}(X)\bigr)
\;\simeq\;
D^b\bigl(\mathrm{Fuk}(X^\vee)\bigr)
\]
for a mirror pair $(X, X^\vee)$.  In this interpretation, the bar-cobar
adjunction should be the passage from B-branes to A-branes: the bar construction
$\barB(\mathcal{A})$ computes the Ext-algebra of the diagonal brane,
and the cobar construction recovers the endomorphism algebra of the dual
collection on $X^\vee$.
\end{conjecture}

\begin{remark}[Scope]
Theorem~\ref*{thm:e1-module-koszul-duality} stands independently; the HMS interpretation requires identifying chiral algebra modules with brane categories (boundary CFT input) and passing from the abstract functor $\Phi$ to derived category equivalences (tilting generation).
\end{remark}

\subsubsection{HMS as Steinberg correspondence duality}
\label{sec:hms-steinberg}
\index{homological mirror symmetry!Steinberg correspondence}
\index{Steinberg correspondence!HMS interpretation}

\providecommand{\BM}{\mathrm{BM}}
\providecommand{\Coh}{\mathrm{Coh}}
\providecommand{\Fuk}{\mathrm{Fuk}}
\providecommand{\HW}{\mathrm{HW}}

The Chriss--Ginzburg perspective on convolution algebras
(cf.\ \S\ref*{sec:steinberg-motivation}) furnishes a geometric
refinement of Conjecture~\ref{conj:hms-koszul}: both sides of HMS
are Steinberg correspondences, and the mirror map is the exchange of
correspondence factors.

\begin{construction}[B-model Steinberg object]
\label{con:steinberg-B}
Let $X$ be a Calabi--Yau manifold with diagonal embedding
$\Delta_X \hookrightarrow X \times X$.  The \emph{B-model Steinberg
object} is the derived self-intersection
\[
  \mathfrak{S}_B
  \;=\;
  \Delta_X \times_{X \times X} \Delta_X
  \;\subset\;
  T^*X \times T^*X.
\]
Its Borel--Moore homology computes the Hochschild cohomology of the
derived category:
\[
  H_*^{\BM}(\mathfrak{S}_B)
  \;=\;
  \mathrm{Ext}^{\bullet}(\mathcal{O}_{\Delta},\,\mathcal{O}_{\Delta})
  \;=\;
  \HH^{\bullet}(\Coh(X)).
\]
The convolution product on $H_*^{\BM}(\mathfrak{S}_B)$ is the
composition of Fourier--Mukai kernels supported on the diagonal.
\end{construction}

\begin{construction}[A-model Steinberg object]
\label{con:steinberg-A}
For the mirror $X^\vee$, fix a Lagrangian $L \subset X^\vee$.
The \emph{A-model Steinberg object} is the Lagrangian
self-intersection
\[
  \mathfrak{S}_A
  \;=\;
  L \times_{X^\vee} L
  \;\subset\;
  T^*X^\vee \times T^*X^\vee.
\]
The wrapped Floer cohomology $\HW^{\bullet}(L, L)$ is the
A-model convolution algebra: the product is the pair-of-pants
cobordism, and the $A_\infty$ structure records all
higher-order disk instantons.
\end{construction}

\begin{proposition}[HMS as Steinberg identification; \ClaimStatusConjectured]
\label{prop:hms-steinberg}
\index{Steinberg correspondence!mirror identification}
Homological mirror symmetry
$D^b(\Coh(X)) \simeq D^b(\Fuk(X^\vee))$
is the identification $\mathfrak{S}_B \simeq \mathfrak{S}_A$
as Steinberg correspondences.  At the level of convolution algebras,
the module bar-cobar functor transports representations:
\[
  \barB\bigl(\HH^{\bullet}(\Coh(X))\bigr)\text{-comodules}
  \;\simeq\;
  \Omega\bigl(\HW^{\bullet}(L, L)\bigr)\text{-modules}.
\]
This is the module Koszul duality statement of
Theorem~\textup{\ref*{thm:e1-module-koszul-duality}},
specialized to the HMS pair $(X, X^\vee)$.
\end{proposition}

\begin{remark}[Swiss-cheese color exchange]
\label{rem:hms-color-exchange}
The bar complex $\barB(\Coh(X))$ computing
$\mathrm{Ext}^{\bullet}(\mathcal{O}_\Delta,\,\mathcal{O}_\Delta)$
\emph{is} the Steinberg object $\mathfrak{S}_B$: the bar differential
encodes the holomorphic convolution, while the bar coproduct encodes
the topological deconcatenation.  Under the mirror map
$X \leftrightarrow X^\vee$, the two factors of the correspondence
exchange roles---holomorphic becomes topological and vice versa.
This is precisely the Swiss-cheese color exchange: the closed
(holomorphic) color on $X$ maps to the open (topological) color
on $X^\vee$, as the Steinberg variety presents the Hecke algebra
and the bar complex presents the Swiss-cheese algebra.
\end{remark}


\subsection{AGT correspondence connections}
\label{sec:agt}

\begin{theorem}[AGT correspondence; \ClaimStatusProvedElsewhere{} \cite{AGT09}]\label{thm:agt}
(Alday--Gaiotto--Tachikawa) For 4d $\mathcal{N}=2$ $SU(2)$ gauge theory with $N_f=4$: $Z_{\mathrm{Nekrasov}}(\epsilon_1,\epsilon_2,a;q) = \langle V_{\alpha_1}(z_1)\cdots V_{\alpha_4}(z_4)\rangle_{\mathrm{Liouville}}$, with $c = 1+6Q^2$, $Q=b+b^{-1}$, $b^2=\epsilon_1/\epsilon_2$.
\end{theorem}

\begin{corollary}[Virasoro from gauge theory; \ClaimStatusProvedElsewhere{} \cite{AGT09}]\label{cor:virasoro-gauge}
The Virasoro algebra ($\Einf$-chiral) arises from the AGT computation; $c$ is determined by $\Omega$-background parameters.
\end{corollary}


\begin{conjecture}[\texorpdfstring{$q$}{q}-AGT; \ClaimStatusConjectured]\label{conj:q-agt}
The 5d lift of AGT (adding a circle) gives:
\[
Z_{\mathrm{5d Nekrasov}}(q, t) = \text{(conformal blocks of } \mathcal{W}_{q,t})
\]
where $\mathcal{W}_{q,t}$ is the quantum W-algebra, an $\Eone$-chiral algebra.
\end{conjecture}

\begin{remark}[Scope]
$\mathcal{W}_{q,t}$ as an $\Eone$-chiral algebra is well-defined (Frenkel--Reshetikhin); the equality with 5d Nekrasov partition functions~\cite{AY10} is a physics conjecture beyond the methods of this monograph.
\end{remark}

\subsection{The 5d Steinberg and quantum toroidal algebras}
\label{sec:5d-steinberg-toroidal}
\index{Steinberg variety!5d instanton}
\index{quantum toroidal algebra}
\index{q-AGT correspondence@$q$-AGT correspondence!Steinberg approach}

\providecommand{\Minst}{\mathcal{M}_{\mathrm{inst}}}

The $q$-AGT conjecture (Conjecture~\ref{conj:q-agt}) admits a
Chriss--Ginzburg--style attack via the 5d instanton Steinberg variety.
We outline the construction and its place in the dimension-cohomology tower.

\begin{construction}[The 5d instanton Steinberg]
\label{constr:5d-steinberg}
\index{Steinberg variety!K-theoretic convolution}
The 5d Nekrasov partition function $Z_{\mathrm{5d}}(q,t)$ is computed by
K-theoretic integration over the instanton moduli space
$\Minst^{\mathrm{5d}}$.  The Hecke correspondence
\[
  \Minst^{\mathrm{5d}}(n)
  \xleftarrow{\;\pi_1\;}
  \mathfrak{H}^{\mathrm{5d}}(n,n+1)
  \xrightarrow{\;\pi_2\;}
  \Minst^{\mathrm{5d}}(n+1)
\]
defines, by fibre product, the 5d instanton Steinberg variety
\[
  \mathfrak{S}^{\mathrm{5d}}
  \;=\;
  \mathfrak{H}^{\mathrm{5d}}
  \times_{\Minst^{\mathrm{5d}}}
  \mathfrak{H}^{\mathrm{5d}}.
\]
The convolution product on $\mathfrak{S}^{\mathrm{5d}}$ is
\emph{K-theoretic}, not homological: the relevant functor is
$(\pi_2)_* \circ \pi_1^*$ acting on K-theory classes, with the
push-forward defined via the equivariant K-theoretic Gysin map.
The convolution algebra is $K_0^{T}(\mathfrak{S}^{\mathrm{5d}})$
with the K-theoretic convolution product.
\end{construction}

\begin{proposition}[K-theoretic convolution and quantum toroidal algebras;
  \ClaimStatusConjectured]
\label{prop:5d-steinberg-toroidal}
\index{Schiffmann--Vasserot}
The K-theoretic convolution algebra
$K_0^{T}(\mathfrak{S}^{\mathrm{5d}})$ is isomorphic to the quantum
toroidal algebra $U_{q,t}(\hat{\hat{\mathfrak{g}}})$ of
Schiffmann--Vasserot.  Under this identification:
\begin{enumerate}
\item The bar complex $\barB(\mathcal{W}_{q,t})$ of the quantum
  W-algebra is the Steinberg object $\mathfrak{S}^{\mathrm{5d}}$:
  bar elements of degree $k$ parametrise $k$-step Hecke chains, and
  the bar differential encodes the K-theoretic convolution product.
\item The $q$-parameter is the K-theory equivariant parameter---the
  exponential of the $\Omega$-background circle radius---replacing
  the homological spectral parameter $z$ of the 4d story.
\item Bar-cobar inversion
  $\Omega\bigl(\barB(\mathcal{W}_{q,t})\bigr) \simeq \mathcal{W}_{q,t}$
  recovers the quantum W-algebra from its Steinberg presentation,
  upgrading the 4d bar-cobar adjunction
  (Theorem~\ref{thm:agt-2d-bar-cobar}) to the K-theoretic setting.
\end{enumerate}
\end{proposition}

\begin{remark}[The dimension-cohomology tower]
\label{rem:dim-cohomology-tower}
\index{dimension-cohomology tower}
\index{elliptic cohomology!Steinberg variety}
The passage from 4d to 5d is the upgrade from Borel--Moore homology to
K-theory on the Steinberg variety.  This is the first step in a
dimension-cohomology tower:
\[
\renewcommand{\arraystretch}{1.4}
\begin{array}{r@{\;:\;}l@{\;\;=\;\;}l@{\qquad}l}
  \text{4d} & H_*^{\mathrm{BM}}(\mathfrak{S}^{\mathrm{4d}})
    & \mathcal{W}\text{-algebra}
    & \text{(homological)} \\
  \text{5d} & K_0(\mathfrak{S}^{\mathrm{5d}})
    & \mathcal{W}_{q,t}\text{-algebra}
    & \text{(K-theoretic)} \\
  \text{6d} & \mathrm{Ell}(\mathfrak{S}^{\mathrm{6d}})
    & \mathcal{W}_{q,p}\text{-algebra}
    & \text{(elliptic cohomology)}
\end{array}
\]
Each row is an instance of the Chriss--Ginzburg paradigm applied
to the instanton moduli space of the corresponding gauge theory:
the convolution algebra of the Steinberg variety, computed in the
appropriate cohomology theory, recovers the symmetry algebra.
The bar complex in each case is the Steinberg object itself; the
cohomology theory determines whether the bar differential and
$\Ainf$-operations are polynomial ($H_*^{\mathrm{BM}}$),
trigonometric ($K_0$), or elliptic ($\mathrm{Ell}$).
\end{remark}

%================================================================
% SECTION: HOLOMORPHIC-TOPOLOGICAL BOUNDARY CONDITIONS
% (from holomorphic_topological.tex)
%================================================================

\section{Holomorphic-topological boundary conditions and 4d origins}
\label{ch:ht-boundary}

\index{holomorphic-topological theory|textbf}

Every chiral algebra in Part~II arises as the boundary theory of a
holomorphic-topological field theory in the sense of Costello--Li.
The bar-cobar adjunction, from this perspective, encodes boundary-bulk
duality: the bar construction computes the open-string state space,
and bar-cobar inversion recovers the boundary algebra from bulk data.

\subsection{Mathematical relationships between frameworks}
\label{sec:precise-relationships-frameworks}

\subsubsection{From 4D gauge theory to 2D chiral algebras}

\begin{theorem}[Costello--Li dimensional reduction; \ClaimStatusProvedElsewhere]\label{thm:costello-li-reduction}
\index{Costello--Li construction}
\cite{CL16} Consider 4D $\mathcal{N}=2$ super Yang--Mills with gauge group $G$
on $\mathbb{C}^2$. After holomorphic-topological twist:
\begin{enumerate}
\item Fields become $\bar{\partial}$-closed differential forms with values in
      $\mathfrak{g} \otimes \mathcal{O}_{\mathbb{C}^2}$

\item The action becomes BV-BRST exact up to a topological term: $S = \{Q_{\text{BRST}}, \Psi\} + S_{\mathrm{top}}$

\item Replacing $\mathbb{C}^2$ by $\mathbb{C} \times \mathbb{C}^*$ and taking $S^1$-equivariant reduction along $|w| = 1 \subset \mathbb{C}^*$
      produces a factorization algebra on $\mathbb{C}$

\item The resulting 2D theory has structure of a \emph{factorization algebra}
      $\mathcal{F}_G$, which is not a priori a chiral algebra
\end{enumerate}
\end{theorem}

\begin{remark}[Factorization algebras versus chiral algebras]\label{rem:fact-vs-chiral}
The distinction is as follows (see \BDref{§2.3, §3.2}):

\emph{Factorization algebra} $\mathcal{F}$: assigns $\mathcal{F}(U)$ to every open set $U \subset X$, with multiplication maps $\mathcal{F}(U) \otimes \mathcal{F}(V) \to \mathcal{F}(U \sqcup V)$ for disjoint $U, V$ satisfying associativity (the factorization property over multiple disjoint opens).  Example: observables in any QFT.

\emph{Chiral algebra} $\mathcal{A}$: assigns a $\mathcal{D}_X$-module $\mathcal{A}_x$ at each point $x \in X$, with chiral product $j_*j^*(\mathcal{A} \boxtimes \mathcal{A}) \to \Delta_!\mathcal{A}$ having prescribed pole structure along the diagonal.  Example: vertex algebras, affine Kac--Moody algebras (Virasoro action is additional structure, not part of the definition).

\emph{Relationship} \cite{BD04,CG17}.
\[\text{Chiral algebras} \xleftrightarrow{\sim} \text{Factorization algebras on curves}\]
is an equivalence (Beilinson--Drinfeld \cite[Theorem~3.4.9]{BD04}).  In higher dimensions, chiral algebras embed into CG-style factorization algebras as those with $\mathcal{D}$-module structure and holomorphic dependence.
\end{remark}

\begin{conjecture}[When does CL produce chiral algebras?; \ClaimStatusConjectured]\label{conj:CL-produces-chiral}
The Costello--Li construction produces a \emph{genuine chiral algebra} (not just
factorization algebra) if and only if:
\begin{enumerate}
\item The 4D theory has additional supersymmetry ensuring holomorphicity
\item The dimensional reduction preserves conformal symmetry
\item Central charge and anomaly terms satisfy consistency conditions
\end{enumerate}

For $\mathcal{N}=4$ SYM, this is established:
Theorem~\ref{thm:cl-n4-chirality} proves that
$\mathcal{F}_{G,\mathrm{Ad}}$ is genuinely chiral, producing the
affine Kac--Moody algebra $\widehat{\mathfrak{g}}_k$ at level~$k$
determined by the gauge coupling.  For $\mathcal{N}=2$ SYM with
general matter content, the precise conditions are conjectured in
Conjecture~\ref{conj:cl-general-chirality}.
\end{conjecture}

\begin{remark}[Evidence]
Chirality requires: (1) the twist preserves a holomorphic structure on the Coulomb branch \cite{Gai19}; (2) $T(z)$ survives the twist with vanishing anomalies; (3) the factorization algebra extends to a $\mathcal{D}$-module on $\mathrm{Ran}(X)$ (automatic for chiral algebras by \BDref{Chapter 3}, requires verification for twisted gauge theories).
\end{remark}

\begin{remark}[Scope]
Chirality is proved for $\mathcal{N}=4$ SYM (Theorem~\ref{thm:cl-n4-chirality}); the general case remains conjectural.  See \cite{CDG20,GKW24,Zeng23} for related work.  MC5 is proved at all genera in Volume~I; the remaining obstruction is the chirality verification for general twisted gauge theories.
\end{remark}

\subsubsection{Costello--Li conditions: from factorization to chiral}

\begin{theorem}[CL chirality for \texorpdfstring{$\mathcal{N}=4$}{N=4}; \ClaimStatusProvedElsewhere]
\label{thm:cl-n4-chirality}
\index{Costello--Li construction!chirality criterion}
\cite{CL16}
Let $G$ be a reductive group and consider 4d $\mathcal{N}=2$ gauge theory
with gauge group $G$ and matter in a representation $R$ on $\Sigma \times \mathbb{C}$.
After holomorphic twist, the theory produces a factorization algebra
$\mathcal{F}_{G,R}$ on $\Sigma$.  When the theory has $\mathcal{N}=4$
supersymmetry (equivalently, $R = \mathrm{Ad}(G)$), the factorization algebra
is strictly chiral: $\mathcal{F}_{G,\mathrm{Ad}}$ satisfies the chiral algebra
axioms of Beilinson--Drinfeld.  The resulting chiral algebra is the affine
Kac--Moody algebra $\widehat{\fg}_k$ at level $k$ determined by the gauge coupling.
\end{theorem}

\begin{conjecture}[CL chirality conditions for general \texorpdfstring{$R$}{R}; \ClaimStatusConjectured]
\label{conj:cl-general-chirality}
\index{Costello--Li construction!$\mathcal{N}=2$ conditions}
For $\mathcal{N}=2$ gauge theory with gauge group $G$ and matter content $R$,
the factorization algebra $\mathcal{F}_{G,R}$ is chiral provided:
\begin{enumerate}[label=\textup{(\roman*)}]
\item $R$ is self-dual: $R \cong R^*$ as a $G$-representation.
\item The one-loop anomaly $c_1(R) \in H^2(\Sigma, \mathbb{Z})$ vanishes.
\item The matter content is balanced: $\chi(\Sigma, R) = 0$ for all genera.
\end{enumerate}
When these conditions hold, the resulting chiral algebra $\mathcal{A}_{\mathrm{CL}}(G, R)$
on $\Sigma$ has:
\begin{itemize}
\item Central charge $c = \dim(G) \cdot (1 - h^\vee/k) + \tfrac{1}{2}\dim(R)$,
  where $k$ is the level and $h^\vee$ is the dual Coxeter number.
\item Koszul dual $\mathcal{A}_{\mathrm{CL}}(G, R)^! = \mathcal{A}_{\mathrm{CL}}(G^\vee, R^\vee)$,
  where $G^\vee$ is the Langlands dual group and $R^\vee$ the dual representation.
\item Bar complex $\barB(\mathcal{A}_{\mathrm{CL}})$ computing the derived moduli
  of solutions to the BPS equations on $\Sigma \times \mathbb{C}$.
\end{itemize}
For $G = \mathrm{GL}(N)$ and $R = \mathrm{Ad}$, the chiral algebra is the
$\mathcal{W}_N$-algebra, and the instanton partition function on $\mathbb{C}^2$
equals the genus-$0$ bar complex evaluation --- this is the AGT correspondence
(Theorem~\ref{thm:agt}).
\end{conjecture}

\begin{remark}[Scope]
Conditions~(i)--(iii) are established for $\mathcal{N}=4$; the general case is conjectural.
\end{remark}

\subsubsection{Paquette--Williams boundary vertex algebras}

\begin{theorem}[Paquette--Williams 2022; \ClaimStatusProvedElsewhere]\label{thm:paquette-williams-boundaries}
\cite{PW22} Consider holomorphic-topological 4D $\mathcal{N}=2$ gauge theory on
$\mathbb{C}^2$ with boundary at $z_2 = 0$. Then:
\begin{enumerate}
\item Boundary conditions $\mathcal{B}$ correspond to Lagrangian submanifolds in
      Coulomb branch $\mathcal{M}_C$

\item The boundary supports a \emph{vertex operator algebra} (VOA) $V_{\mathcal{B}}$

\item This VOA is the quantization of the symplectic reduction of $\mathcal{M}_C$
      at the boundary

\item The VOA $V_{\mathcal{B}}$ has a \emph{chiral envelope} $\mathcal{A}_{\mathcal{B}}$,
      which is a chiral algebra in the BD sense
\end{enumerate}
\end{theorem}

\begin{remark}[Bar-cobar duality for boundary VOAs]\label{rem:PW-connection}
Paquette--Williams produce vertex algebras, which are algebraic objects (Frenkel--Ben-Zvi
\cite{FBZ04}). The bar-cobar framework of Part~I applies to their \emph{chiral envelopes}, which are
geometric objects ($\mathcal{D}$-modules).

\begin{center}
\begin{tikzcd}
\text{Boundary VOA } V_{\mathcal{B}} \ar[r, "\text{chiral envelope}"] \ar[d, "\text{PW}"]
& \text{Chiral algebra } \mathcal{A}_{\mathcal{B}} \ar[d, "\text{bar-cobar}"] \\
\text{4D HT gauge theory} \ar[r, "\text{CL reduction}"]
& \text{2D factorization algebra}
\end{tikzcd}
\end{center}

The vertex algebra $V_{\mathcal{B}}$ contains the algebraic information (OPE, modes,
etc.), while the chiral algebra $\mathcal{A}_{\mathcal{B}}$ contains the geometric
information (configuration spaces, $\mathcal{D}$-modules, sheaf cohomology).

Bar-cobar duality (Part~I) applies to $\mathcal{A}_{\mathcal{B}}$: the bar construction $\bar{B}(\mathcal{A}_{\mathcal{B}})$ \emph{is} the factorization coalgebra on $\Ran(X)$ that encodes PW boundary data (Convention~\ref*{conv:bar-coalgebra-identity}).  Verdier duality $\mathbb{D}_{\Ran}$ converts this coalgebra into a factorization algebra; on the Koszul locus, the cohomology of $\mathbb{D}_{\Ran}\bar{B}(\mathcal{A}_{\mathcal{B}})$ is the Koszul dual $\mathcal{A}_{\mathcal{B}}^!$, identifying the dual boundary condition.  Separately, the cobar--bar counit $\Omega\bar{B}(\mathcal{A}_{\mathcal{B}}) \xrightarrow{\sim} \mathcal{A}_{\mathcal{B}}$ recovers the original algebra (this is bar-cobar inversion, a different operation from Verdier duality).
\end{remark}

\subsection{From 4d SYM to holomorphic Chern--Simons}

\subsubsection{The holomorphic twist and localization}

\begin{definition}[Holomorphic twist of 4d \texorpdfstring{$\mathcal{N}=4$}{N=4} SYM]
Start with $\mathcal{N}=4$ super Yang--Mills in 4d with gauge group $G$, viewed as an $\mathcal{N}=2$ theory by choosing a preferred $\mathcal{N}=2$ subalgebra (Costello~\cite{costello-renormalization}).  The field content before twisting:
\begin{itemize}
\item Vector multiplet: $A_\mu$ (gauge field), $\Phi^I$ (6 scalars), $\lambda,
\bar{\lambda}$ (fermions)
\item Supersymmetry: 16 supercharges transforming under $\text{Spin}(6)_R$
\end{itemize}

The \emph{holomorphic twist} (not to be confused with the fully topological Kapustin--Witten A-twist of $\mathcal{N}=4$):
\begin{enumerate}
\item Decompose $\mathbb{R}^4 = \mathbb{C} \times \mathbb{C}$ with coordinates
$(z, w)$
\item Twist the Lorentz group: $\text{SO}(4) \to \text{SO}(2)_{\text{hol}} \times
\text{SO}(2)_{\text{top}}$
\item Mix with R-symmetry to make a supercharge $Q$ scalar
\item Result: Theory is holomorphic in $z$, topological in $w$
\end{enumerate}
\end{definition}

\begin{theorem}[Localization to holomorphic data; \ClaimStatusProvedElsewhere]\label{thm:ht-localization}
After the holomorphic twist \cite{costello-renormalization}, the path integral localizes to:
\[Z = \int_{[\bar{\partial}_A = 0]} \mathcal{O}(\text{fields}) \cdot e^{-S_{\text{inst}}}\]

where $\bar{\partial}_A = 0$ means:
\begin{itemize}
\item $A$ is a holomorphic connection
\item Scalars $\Phi$ satisfy holomorphic moment map equation
\end{itemize}

This is the moduli space of holomorphic $G$-bundles with Higgs field.
\end{theorem}

\begin{proof}[Sketch]
The twisted action decomposes as $S_{\mathrm{twisted}} = \{Q, V\} + S_0$ where $Q$ is the scalar supercharge, $V$ is the gauge fermion, and $S_0$ is metric-independent.  Since $\{Q, V\} \geq 0$, the path integral localizes to
\[\mathcal{M}_Q = \{F_{A}^{(0,2)} = 0,\; F_A^{(1,1)} + [\Phi, \Phi^*] = 0,\; \bar{\partial}_A \Phi = 0\},\]
i.e., Hitchin's self-duality equations in the holomorphic gauge.
\end{proof}

\subsubsection{Holomorphic Chern--Simons as effective theory}

\begin{definition}[Holomorphic Chern--Simons action]
\index{holomorphic Chern--Simons}
On a Calabi--Yau 3-fold $X$ with holomorphic volume form $\Omega \in \Omega^{3,0}(X)$, the
holomorphic Chern--Simons action is:
\[S_{\text{HCS}}[A] = \frac{1}{2}\int_{X} \Omega \wedge \operatorname{Tr}\left(A
\wedge \bar{\partial} A + \frac{2}{3} A \wedge A \wedge A\right)\]

where $A \in \Omega^{0,1}(X, \mathfrak{g}_{\mathbb{C}})$.  Since $A \in \Omega^{0,1}$ and $\bar\partial A \in \Omega^{0,2}$, we have $A \wedge \bar\partial A \in \Omega^{0,3}$ and $A^3 \in \Omega^{0,3}$, so $\Omega \wedge \text{Tr}(\cdots)$ is a $(3,3)$-form, hence integrable over~$X$.  The factor $\frac{1}{2}$ follows the convention of Costello~\cite{costello-renormalization}.
\end{definition}

\begin{theorem}[HCS from dimensional reduction; \ClaimStatusProvedElsewhere]\label{thm:hcs-dimensional-reduction}
Holomorphic Chern--Simons \cite{costello-renormalization} arises from the holomorphic twist of 4d SYM by:
\begin{enumerate}
\item The holomorphic twist on $\mathbb{C}^2$ yields a theory holomorphic in $z$, topological in~$w$
\item Compactify the topological direction $w$ on $S^1$ to produce a 3d theory on $\mathbb{C} \times \mathbb{R}$
\item This 3d theory is holomorphic Chern--Simons
\end{enumerate}

The 4d holomorphic twist yields a holomorphic-topological theory on $\mathbb{C}^2$. Compactifying one direction produces a 3d theory; the holomorphic volume form on the resulting CY 3-fold $X$ is a $(3,0)$-form $\Omega \in \Omega^{3,0}(X)$.  Locally, taking $X = \mathbb{C}^3$ with coordinates $(z, w, u)$:
\[\Omega = dz \wedge dw \wedge du\]
\end{theorem}

\subsection{Boundary conditions and chiral operads}

\subsubsection{The deformed conifold geometry}

\begin{example}[Costello--Gaiotto holography setup]
The canonical example is the deformed conifold:
\[X = \{u_1 w_2 - u_2 w_1 = N\} \subset \mathbb{C}^4\]

This space has:
\begin{itemize}
\item Calabi--Yau 3-form: $\Omega|_X = \mathrm{Res}_{f=N}\frac{du_1 \wedge du_2 \wedge dw_1
\wedge dw_2}{f - N}$ (Poincar\'e residue of the ambient 4-form along $X = \{f = N\}$, yielding a holomorphic 3-form on $X$)
\item $SL_2(\mathbb{C})$ isometry: Acting on $(u, w)$ coordinates
\item Asymptotic boundary: $\mathbb{CP}^1 \times \mathbb{CP}^1$ as $|u|, |w|
\to \infty$
\item Pole structure: $\Omega$ has cubic pole along the boundary divisor
\end{itemize}
\end{example}


\subsubsection{HT boundary conditions}

\begin{definition}[Holomorphic-topological boundary condition]
For HCS on $X$ with boundary compactification $\bar{X}$ and boundary divisor
$D = \bar{X} \setminus X$:

A \emph{holomorphic-topological boundary condition} specifies:
\begin{enumerate}
\item Extension: Fields extend to $\bar{X}$ as holomorphic sections
\item Vanishing order: $A \in H^0(\bar{X}, \Omega^{0,1}(\mathcal{O}(-D)))$ (simple zero at $D$)
\item Behavior: As approaching $D$, $A \sim (\text{distance}) \cdot (\text{smooth})$
\end{enumerate}
\end{definition}

\begin{theorem}[Boundary chiral algebra; \ClaimStatusProvedElsewhere]
\label{thm:ht-boundary-chiral-algebra}
\cite{CDG20}
An HT boundary condition supports a chiral algebra $\mathcal{A}_{\text{bdy}}$ whose:
\begin{enumerate}
\item \emph{Generators.} Boundary local operators $\mathcal{O}(z)$ for $z \in D$
\item \emph{OPE.} Determined by bulk path integral with boundary insertions
\[\mathcal{O}_1(z) \mathcal{O}_2(w) = \sum_k C_{12}^k(z-w) \mathcal{O}_k(w)\]
\item \emph{Factorization.} Extends to factorization algebra on $D$
\end{enumerate}
\end{theorem}

\begin{proof}[Sketch]
Boundary operators $\mathcal{O}(z) = \lim_{\epsilon \to 0} \mathrm{Tr}(A(z + \epsilon n) \cdots)$ are well-defined by the simple zero condition.  The OPE arises from short-distance singularities of the HCS path integral with boundary insertions.  Locality (annular convergence), associativity (path integral composition), and skyscraper support on~$D$ give the BD chiral algebra axioms.
\end{proof}

\subsubsection{Chiral operad action}

\begin{conjecture}[Chiral operad from HCS; \ClaimStatusConjectured]
\label{conj:chiral-operad-hcs}
The holomorphic Chern--Simons theory defines a chiral operad $\mathcal{P}_{\text{HCS}}$
acting on boundary chiral algebras.

The operad operations:
\[\mathcal{P}_{\text{HCS}}(n) = \text{Obs}(X; D \times \overline{C}_n(D))\]

are observables on $X$ with marked points on the boundary.
\end{conjecture}

\begin{remark}[Scope]
The operad structure is predicted by \cite{CG17}; the explicit identification with HCS observables requires perturbative quantization.
\end{remark}

\begin{example}[Kac--Moody from gauge HCS]
For $G = SU(N)$ holomorphic Chern--Simons:
\[\mathcal{A}_{\text{bdy}} = \widehat{\mathfrak{sl}}_N\]

the affine Kac--Moody algebra at level $k$ (determined by HCS coupling).

The boundary currents:
\[J^a(z) = \lim_{\epsilon \to 0} \text{Tr}(T^a A(z + \epsilon n))\]

satisfy the Kac--Moody OPE:
\[J^a(z) J^b(w) \sim \frac{k \delta^{ab}}{(z-w)^2} +
\frac{f^{abc} J^c(w)}{z-w}\]

This is derivable from the HCS path integral.
\end{example}

\subsection{Open-closed correspondence as bar-cobar duality}

\subsubsection{Open and closed sectors}

\begin{theorem}[Open-string bar identification; \ClaimStatusProvedHere]
\label{thm:open-string-bar}
\index{open-closed correspondence!open sector}
Let $\mathcal{A}_{\mathrm{bdy}}$ be the boundary chiral algebra of
a holomorphic-topological theory with boundary condition supporting
a Kac--Moody or W-algebra vertex algebra.  Then the bar complex
$\barB^{\mathrm{ch}}(\mathcal{A}_{\mathrm{bdy}})$ computes the
open-string state space: there is a quasi-isomorphism
\[
\barB^{\mathrm{ch}}(\mathcal{A}_{\mathrm{bdy}})
\simeq
C^{\frac{\infty}{2}+\bullet}(\mathcal{A}_{\mathrm{bdy}})
\]
to the semi-infinite (BRST) complex of~$\mathcal{A}_{\mathrm{bdy}}$.
Moreover, Verdier duality on $\operatorname{Ran}(X)$ converts the
bar coalgebra $\barB^{\mathrm{ch}}(\cA_{\mathrm{bdy}})$ into a
factorization algebra whose cohomology, on the Koszul locus, is the
Koszul dual~$\cA_{\mathrm{bdy}}^!$
\textup{(}Theorem~\textup{\ref*{thm:verdier-bar-cobar})}.
The coalgebra encodes boundary data; its Verdier dual algebra
encodes the bulk.
\end{theorem}

\begin{proof}
For Kac--Moody boundary conditions
$\mathcal{A}_{\mathrm{bdy}} = \widehat{\fg}_k$ ($k \neq -h^\vee$),
Theorem~\ref*{thm:bar-semi-infinite-km} provides the quasi-isomorphism
$\barB^{\mathrm{ch}}(\widehat{\fg}_k) \simeq C^{\frac{\infty}{2}+\bullet}(\widehat{\fg}_k)$
via the PBW filtration and Arkhipov's functor.
For W-algebra boundary conditions
$\mathcal{A}_{\mathrm{bdy}} = \mathcal{W}_k(\fg)$,
Theorem~\ref*{thm:bar-semi-infinite-w} gives the analogous
identification through Drinfeld--Sokolov reduction.
In both cases, the semi-infinite complex is the mathematical
formulation of the open-string BRST complex: its cohomology
computes the physical state space of open strings ending on
the holomorphic-topological boundary condition.

For the Verdier duality statement:
Theorem~\ref*{thm:verdier-bar-cobar} shows that $\mathbb{D}_{\operatorname{Ran}}$
maps the bar coalgebra to a factorization algebra whose cohomology
on the Koszul locus is $\cA_{\mathrm{bdy}}^!$.
Verdier duality converts the coalgebra to an algebra; this is
distinct from the cobar functor $\Omega$ (the left adjoint to~$\barB$,
which recovers~$\cA$ by bar-cobar inversion).
\end{proof}

\begin{conjecture}[Closed-string cobar identification; \ClaimStatusConjectured]
\label{conj:closed-string-cobar}
\index{open-closed correspondence!closed sector}
Let $\mathcal{C}_{\mathrm{bulk}}$ be the bulk chiral coalgebra of
a holomorphic-topological string theory.  Then:
\begin{enumerate}
\item \emph{Closed strings.} The cobar complex
  $\Omega^{\mathrm{ch}}(\mathcal{C}_{\mathrm{bulk}})$ computes the
  closed-string field theory state space in the sense of
  Zwiebach~\cite{Zwi93}.
\item \emph{Open-closed duality.} The bar-cobar adjunction
  should furnish the full open-closed correspondence: the natural map
  $\barB(\mathcal{A}_{\mathrm{bdy}}) \to
  \Omega(\mathcal{C}_{\mathrm{bulk}})$ is conjecturally the open-closed string
  field theory map.
\end{enumerate}
\end{conjecture}

\begin{remark}[Evidence]
Open sector: boundary vertex operators are elements of $\mathcal{A}_{\mathrm{bdy}}$, and off-shell amplitudes lie in $\barB^{\mathrm{ch}}(\mathcal{A}_{\mathrm{bdy}})$ (proved).  Closed sector: bulk vertex operators give distribution-valued correlation functions on open configuration spaces (conjectural).  The open-closed map $\barB(\mathcal{A}_{\mathrm{bdy}}) \to \Omega(\mathcal{C}_{\mathrm{bulk}})$ sends boundary residues to bulk distributions via Stokes' theorem.
\end{remark}

\begin{remark}[Scope]
Open-string side proved (Theorems~\ref*{thm:bar-semi-infinite-km}, \ref*{thm:bar-semi-infinite-w}).  MC4 and MC5 are both proved in Volume~I; the closed-string identification is now accessible via the all-genera bar-cobar framework.
\end{remark}

\subsubsection{Factorization and dimensional reduction}

\begin{conjecture}[Factorization along dimension tower; \ClaimStatusConjectured]\label{conj:factorization-dim-tower}
The dimensional reduction sequence:
\[\text{4d SYM} \xrightarrow{\text{HT twist + reduce}}
\text{3d HCS} \xrightarrow{\text{boundary}}
\text{2d chiral algebra} \xrightarrow{\text{defect}}
\text{1d quantum mechanics}\]

is governed by iterated bar-cobar constructions at each level.
\end{conjecture}

\begin{remark}[Scope]
Individual levels are established \cite{BD04,CG17,CWY18}; MC5 is proved at all genera in Volume~I, so the bar-cobar side of the iterated identification is now unconditional.
\end{remark}

\begin{example}[Explicit tower for \texorpdfstring{$\mathcal{N}=4$}{N=4} SYM]
$\mathcal{N}=4$ SYM with gauge group $G$ on $\mathbb{R}^4$
$\xrightarrow{\text{HT twist + reduce}}$
HCS on $\mathbb{C} \times \mathbb{R}_{\geq 0}$
$\xrightarrow{\text{boundary}}$
$\widehat{\mathfrak{g}}_k$ on $\mathbb{C}$
$\xrightarrow{\text{defect}}$
quantum integrable system (Toda for $G = SU(N)$).
\end{example}

\subsubsection{The iterated Steinberg and the dimension tower}

The conjecture above admits a precise operadic formulation via the
Chriss--Ginzburg convolution paradigm, applied level by level to the
dimension tower.

\providecommand{\Steinberg}{\mathfrak{S}}

\begin{construction}[Steinberg object at level $k$]
\label{constr:steinberg-level-k}
For an HT theory on $\mathbb{C}^k \times \mathbb{R}$, define the
\emph{level-$k$ Steinberg object} $\Steinberg_b^{(k)}$ as the
convolution algebra whose correspondence space at arity $n$ is
\[
  \Steinberg_b^{(k)}(n) \;=\; \operatorname{FM}_n(\mathbb{C}^k)
  \times \operatorname{Conf}_n(\mathbb{R}).
\]
The differential comes from OPE residues on $\operatorname{FM}_n(\mathbb{C}^k)$
(holomorphic directions), and the coproduct from ordered deconcatenation on
$\operatorname{Conf}_n(\mathbb{R})$ (topological direction).
The inclusions $\mathbb{C}^k \hookrightarrow \mathbb{C}^{k+1}$
induce \emph{stabilization maps}
\[
  \sigma_k \colon \Steinberg_b^{(k)} \longrightarrow \Steinberg_b^{(k+1)}
\]
via the operadic maps
$\operatorname{FM}_n(\mathbb{C}^k) \hookrightarrow
\operatorname{FM}_n(\mathbb{C}^{k+1})$.
The \emph{dimension tower} is the directed system
$\Steinberg_b^{(1)} \to \Steinberg_b^{(2)} \to \cdots$.
\end{construction}

\begin{proposition}[\ClaimStatusConjectured]
\label{prop:steinberg-tower-factorization}
The bar-cobar pairing factorizes along the dimension tower:
at level~$k$, the Fulton--MacPherson operad
$\operatorname{FM}(\mathbb{C}^k)$ is formal by Kontsevich formality,
with homology the operad $\mathcal{P}_{k\text{-}1}$ of
$(k{-}1)$-Poisson algebras.  Consequently the level-$k$ Steinberg
$\Steinberg_b^{(k)}$ computes the $(k{-}1)$-shifted Poisson version
of the bar complex, and the stabilization $\sigma_k$ corresponds to
the canonical map $\mathcal{P}_{k-1} \to \mathcal{P}_k$ of Poisson
operads.
\end{proposition}

\begin{remark}[The first three levels]
\label{rem:three-levels}
The tower at low $k$ recovers familiar structures:
\begin{itemize}
\item $k = 1$: $\operatorname{FM}_n(\mathbb{C}^1) \simeq
  \operatorname{FM}_n(\mathbb{C})$.  The homology operad is
  $\mathcal{P}_0 = E_1$, the associative operad.  The Steinberg
  $\Steinberg_b^{(1)}$ is the ordinary associative bar complex.
\item $k = 2$: The homology operad is $\mathcal{P}_1$, i.e.\
  $1$-Poisson $= E_2$.  The Steinberg $\Steinberg_b^{(2)}$ is the
  chiral bar complex of Volume~I---the Swiss-cheese algebra with one
  holomorphic direction $\mathbb{C}$ and one topological direction
  $\mathbb{R}$.
\item $k = 3$: The homology operad is $\mathcal{P}_2$
  ($2$-Poisson).  The Steinberg $\Steinberg_b^{(3)}$ is the
  \emph{2-chiral bar complex}, governing 4d HT theories on
  $\mathbb{C}^3 \times \mathbb{R}$---each level adds one holomorphic
  direction to the Steinberg correspondence.
\end{itemize}
\end{remark}

\subsection{$\mathcal{W}$-algebras from Hitchin moduli}

\subsubsection{The Higgs branch and Hitchin system}

\begin{definition}[Hitchin moduli space]
For a Riemann surface $\Sigma_g$ of genus $g$ and gauge group $G$, the Hitchin
moduli space is:
\[\mathcal{M}_{\text{Hit}}(\Sigma_g, G) = \{(E, \Phi) : \bar{\partial}_E \Phi = 0\} /
\sim\]

where $E \to \Sigma_g$ is a holomorphic $G$-bundle, $\Phi \in H^0(\Sigma_g, \operatorname{End}(E) \otimes K_{\Sigma_g})$ is the Higgs field, and $\sim$ denotes gauge equivalence.
\end{definition}

\begin{proposition}[Virasoro from \texorpdfstring{$SL_2$}{SL_2}-Hitchin; \ClaimStatusProvedElsewhere]\label{prop:virasoro-from-hitchin}
For $G = SL_2$, the chiral algebra of local operators on the Hitchin moduli space $\mathcal{M}_{\mathrm{Hit}}(\Sigma_g, SL_2)$ is the Virasoro algebra:
\[\mathcal{A}_{\mathrm{local}}(\mathcal{M}_{\mathrm{Hit}}(\Sigma_g, SL_2)) \cong \mathrm{Vir}_c.\]
\end{proposition}

\begin{proof}
The Hitchin system for $SL_2$ is the cotangent bundle $T^*\mathrm{Bun}_{SL_2}(\Sigma_g)$; quantization of the coordinate ring via the Beilinson--Drinfeld construction \cite{BD04} produces the Virasoro vertex algebra at the level determined by the symplectic structure.
\end{proof}

\begin{conjecture}[W-algebra from Hitchin; \ClaimStatusConjectured]
\label{conj:w-algebra-hitchin}
For general $G$, the chiral algebra of local operators on $\mathcal{M}_{\mathrm{Hit}}(\Sigma_g, G)$ is
\[
\mathcal{A}_{\mathrm{local}}(\mathcal{M}_{\mathrm{Hit}}) \cong \mathcal{W}(G),
\]
the W-algebra associated to~$G$.  For $G = SL_N$, this is the
$\mathcal{W}_N$ algebra with generators
$T(z)$, $W^{(3)}(z)$, \ldots, $W^{(N)}(z)$
of conformal weights $2, 3, \ldots, N$.
\end{conjecture}

\begin{remark}[Scope]
The $SL_2$ case is Beilinson--Drinfeld \cite{BD04}; the general case requires AGT localization.  Partial results: Schiffmann--Vasserot \cite{SV13}, Maulik--Okounkov \cite{MO19}, Arakawa \cite{Ara07}.
\end{remark}

\begin{proof}[Via AGT correspondence]
Place 4d $\mathcal{N}=2$ gauge theory with gauge group $G$ on $\mathbb{R}^2_\epsilon \times \Sigma_g$ with $\Omega$-background parameters $(\epsilon_1, \epsilon_2)$.  The $\Omega$-background localizes the path integral to $\mathcal{M}_{\mathrm{Hit}}$.  In the Nekrasov limit $\epsilon_2 \to 0$:
\[Z_{4d}|_{\epsilon_2 \to 0} = Z_{\mathcal{W}}[\Sigma_g],\]
the $\mathcal{W}(G)$ partition function on $\Sigma_g$.  The operator correspondence (4d line operators $\leftrightarrow$ 2d vertex operators \cite{DNP25}, 4d surface operators $\leftrightarrow$ 2d extended operators) identifies local operators on $\mathcal{M}_{\mathrm{Hit}}$ with generators of $\mathcal{W}(G)$.
\end{proof}

\subsubsection{Bar-cobar for $\mathcal{W}$-algebras}

\begin{theorem}[\texorpdfstring{$\mathcal{W}$}{W}-algebra bar complex; \ClaimStatusProvedHere]
\label{thm:w-algebra-bar-complex}
For $\mathcal{W}_N$, the geometric bar complex:
\[\bar{B}^{\text{ch}}(\mathcal{W}_N) =
\bigoplus_{n \geq 0} \Gamma\!\left(\overline{C}_{n+1}(X),\,
  j_*j^*\mathcal{W}_N^{\boxtimes(n+1)} \otimes \Omega^n_{\log D}\right)\]

has three proved outputs:
\begin{enumerate}
\item \textup{[\ClaimStatusProvedHere]}
\emph{Conformal blocks.}
The bar-complex cohomology
$H^0(\bar{B}^{\mathrm{ch}}(\mathcal{W}_N))$ computes conformal blocks.
This is the Part~I bar-complex/chiral-homology correspondence applied
to the specific chiral algebra $\mathcal{W}_N$.

\item \textup{[\ClaimStatusProvedHere]}
\emph{Fusion rules.} For rational $\mathcal{W}_N$
\textup{(}$C_2$-cofinite, Arakawa \cite{Ara07}\textup{)},
the bar differential with module insertions at marked
points encodes fusion: $\mathrm{Ext}^1_{\mathcal{W}_N}(M_i
\otimes M_j, M_k)$ is computed by the bar complex
$\barB(\mathcal{W}_N; M_i, M_j, M_k)$, and fusion
multiplicities $N_{ij}^k = \dim
\mathrm{Hom}(M_k, M_i \boxtimes M_j)$ are extracted via
residues at collision divisors.

\item \textup{[\ClaimStatusProvedHere]}
\emph{Modular functors.} The genus-$g$ bar complex
$\barB^{(g)}(\mathcal{W}_N)$ defines a vector bundle
over $\overline{\mathcal{M}}_{g,n}$ whose fibers are
conformal block spaces; the modular operad structure
\textup{(}Theorem~\textup{\ref*{thm:genus-induction-strict})}
governs the genus dependence, and
$\barB^{(g)}$ computes chiral homology
\textup{(}Theorem~\textup{\ref*{thm:genus-g-chiral-homology})}.
\end{enumerate}
\end{theorem}

\begin{remark}[Scope]
All three items are proved.  Item~(2) uses the Huang--Lepowsky--Zhang tensor product theory \cite{HLZ} for $C_2$-cofinite algebras.  Explicit Virasoro modular functor computations are in Theorem~\ref*{thm:chain-modular-functor}.
\end{remark}

\begin{example}[Virasoro = \texorpdfstring{$\mathcal{W}_2$}{W_2}]
For $G = SL_2$, $\mathcal{W}_2 = \text{Vir}$ is the Virasoro algebra.

The bar complex element:
\[\omega = T(z_1) \otimes T(z_2) \otimes \cdots \otimes T(z_n) \otimes
\bigwedge_{i<j} \eta_{ij}^{k_{ij}}\]

represents an off-shell correlator of stress tensors.

The bar differential:
\begin{itemize}
\item $d_{\text{res}}$: Extracts OPE $T(z_1)T(z_2) = \frac{c/2}{(z_1-z_2)^4} + \cdots$
\item $d_{\text{strat}}$: Accounts for degeneration of moduli space
\item $d_{\text{int}}$: Implements Ward identities
\end{itemize}

On-shell correlators (physical observables) are in:
\[H^0(\bar{B}^{\text{ch}}(\text{Vir}))\]
\end{example}

\subsection{The spectral curve Steinberg and $\mathcal{W}$-algebra duality}

\providecommand{\fStein}{\mathfrak{S}}

The Hitchin conjecture (Conjecture~\ref{conj:w-algebra-hitchin}) identifies
$\mathcal{W}(G)$ as the chiral algebra of the Hitchin moduli space, while
Theorem~\ref{thm:w-algebra-bar-complex} computes $\mathcal{W}$-algebra
conformal blocks, fusion rules, and modular functors via the bar complex.
The two results meet in a single geometric object: the \emph{Hitchin
Steinberg variety}, whose convolution algebra \emph{is} the
$\mathcal{W}$-algebra and whose geometric operations \emph{are} the bar
and cobar functors.

\subsubsection{The Hitchin Steinberg}

\begin{construction}[Hitchin Steinberg variety]
\label{constr:hitchin-steinberg}
Let $X$ be a smooth projective curve of genus $g$, $G$ a reductive group,
and $\mathcal{M}_{\mathrm{Hit}}(X,G)$ the Hitchin moduli space.  The Hitchin
fibration
\[
h \colon \mathcal{M}_{\mathrm{Hit}}(X,G) \longrightarrow \mathcal{B}
  = \bigoplus_{i=1}^{\mathrm{rk}\,G} H^0(X, K_X^{\otimes d_i})
\]
has generic fiber a Lagrangian abelian variety --- the Jacobian of
the \emph{spectral curve} $\Sigma_b \subset T^*X$ determined by the
characteristic polynomial of the Higgs field $\varphi$.  For
$b \in \mathcal{B}^{\mathrm{reg}}$, the spectral curve
$\Sigma_b \hookrightarrow T^*X$ is a smooth branched cover of degree
$\mathrm{rk}\,G$.

The \emph{Hitchin Steinberg variety} is the fiber product
\[
\fStein_{\mathrm{Hit}}
  \;=\; \Sigma_b \,\times_{T^*X}\, \Sigma_b
  \;\subset\; T^*X \times T^*X,
\]
the self-intersection of the spectral curve inside the cotangent bundle.
\end{construction}

\begin{conjecture}[$\mathcal{W}$-algebra as convolution on $\fStein_{\mathrm{Hit}}$; \ClaimStatusConjectured]
\label{conj:w-steinberg-convolution}
The Borel--Moore homology of the Hitchin Steinberg variety carries a
convolution product, and there is an isomorphism of vertex algebras
\[
H_*^{\mathrm{BM}}(\fStein_{\mathrm{Hit}})
  \;\cong\; \mathcal{W}_k(\mathfrak{g}),
\]
where the level $k$ is determined by the point $b \in \mathcal{B}$.
The convolution product on $\fStein_{\mathrm{Hit}}$ becomes the
normally-ordered product of $\mathcal{W}_k(\mathfrak{g})$; the
intersection pairing becomes the $\lambda$-bracket.
\end{conjecture}

This is the $\mathcal{W}$-algebra analogue of the classical statement
that the Borel--Moore homology of the Steinberg variety
$\widetilde{\mathcal{N}} \times_{\mathfrak{g}} \widetilde{\mathcal{N}}$
presents the group algebra of the Weyl group (Kazhdan--Lusztig) ---
or, in the affine setting, the affine Hecke algebra
(Chriss--Ginzburg~\cite{CG97}).

\subsubsection{Bar-cobar duality as convolution-restriction}

\begin{proposition}[Bar-cobar $=$ convolution-restriction on $\fStein_{\mathrm{Hit}}$; \ClaimStatusConjectured]
\label{prop:bar-cobar-steinberg}
The bar-cobar duality
$\barB(\mathcal{W}_k) \leftrightarrow \Omega(\mathcal{W}_k^!)$ is the
convolution-restriction adjunction on $\fStein_{\mathrm{Hit}}$:
\begin{enumerate}
\item \emph{Bar differential.}  The bar differential
  $d_{\barB}$ is the convolution map on $\fStein_{\mathrm{Hit}}$:
  collision of points on the spectral curve $\Sigma_b \times \Sigma_b$
  along diagonal strata extracts OPE residues.  In coordinates, for
  sheets $\sigma_i, \sigma_j$ meeting along $\Delta \subset T^*X$,
  \[
  d_{\barB}(\sigma_i \otimes \sigma_j)
    = \operatorname{Res}_{\sigma_i \to \sigma_j}
      \bigl[\sigma_i \cdot \sigma_j\bigr],
  \]
  where the residue is taken along the collision divisor
  $\{p_i = p_j\} \subset \Sigma_b \times \Sigma_b$.

\item \emph{Coproduct.}  The bar coproduct $\Delta_{\barB}$ is the
  restriction map: decomposition of $\fStein_{\mathrm{Hit}}$ into
  individual spectral sheets.  For a class
  $\alpha \in H_*^{\mathrm{BM}}(\fStein_{\mathrm{Hit}})$,
  \[
  \Delta_{\barB}(\alpha) = \sum_{I \sqcup J = \{1,\ldots,N\}}
    \alpha|_{\Sigma_I \times \Sigma_J},
  \]
  where $\Sigma_I = \bigcup_{i \in I} \sigma_i$ is the union of sheets
  indexed by~$I$.
\end{enumerate}
\end{proposition}

\subsubsection{Compatibility with Drinfeld--Sokolov reduction}

\begin{proposition}[DS reduction commutes with the Steinberg construction; \ClaimStatusConjectured]
\label{prop:ds-steinberg-commute}
Let $\mathfrak{n} \subset \mathfrak{g}$ be the nilpotent radical
associated to a nilpotent orbit $e \in \mathfrak{g}$.  Hamiltonian
reduction of $\fStein_{\mathrm{Hit}}$ by the $\mathfrak{n}$-action
produces the reduced Steinberg variety for the DS-reduced
$\mathcal{W}$-algebra.  Explicitly, the following diagram commutes:
\[
\begin{tikzcd}
\fStein_{\mathrm{Hit}}(\mathfrak{g})
  \ar[r, "\text{convolution}"]
  \ar[d, "/\!\!/ \mathfrak{n}"']
& \mathcal{W}_k(\mathfrak{g})
  \ar[d, "\mathrm{DS}_e"] \\
\fStein_{\mathrm{Hit}}(\mathfrak{g})
  /\!\!/ \mathfrak{n}
  \ar[r, "\text{convolution}"']
& \mathcal{W}_k(\mathfrak{g}, e)
\end{tikzcd}
\]
\end{proposition}

\begin{proof}
The $\mathfrak{n}$-action on the spectral curve $\Sigma_b$ is free on the
regular locus $\mathcal{B}^{\mathrm{reg}}$: the spectral curve over a
regular point is smooth, and the $N$-orbits on the fiber
$h^{-1}(b) \cong \operatorname{Jac}(\Sigma_b)$ are free because the
stabilizer of a regular Higgs field under the adjoint $N$-action is
trivial.  Freeness implies that Hamiltonian reduction commutes with the
fiber product:
\[
(\Sigma_b \times_{T^*X} \Sigma_b) /\!\!/ N
  \;\cong\;
(\Sigma_b /\!\!/ N) \times_{T^*X /\!\!/ N} (\Sigma_b /\!\!/ N).
\]
The reduced spectral curve $\Sigma_b /\!\!/ N$ parametrizes the
DS-reduced Hitchin system, so the reduced fiber product is the Steinberg
variety for $\mathcal{W}_k(\mathfrak{g}, e)$.  The geometric
commutativity---that Hamiltonian reduction commutes with fiber
products under the freeness hypothesis---is unconditional.

The horizontal arrows in the diagram invoke
Conjecture~\ref{conj:w-steinberg-convolution} (the identification
of convolution on $\fStein_{\mathrm{Hit}}$ with $\mathcal{W}_k$),
which is why this proposition is conditional.  The algebraic
content---that bar-cobar intertwines with DS reduction at the chain
level---is unconditional via
Theorem~\hyperref[thm:ds-koszul-intertwine]{ds-koszul-intertwine}
of Volume~I.  The present proposition lifts that algebraic statement
to the geometric Steinberg setting, contingent on the Steinberg
presentation conjecture.
\end{proof}

\subsubsection{DS reduction as Hamiltonian reduction of the Steinberg}
\label{subsubsec:ds-hamiltonian-steinberg}

\providecommand{\Lloop}{\mathcal{L}}
\providecommand{\Mvac}{\mathcal{M}_{\mathrm{vac}}}

The preceding proposition establishes that DS reduction commutes with
the Steinberg construction at the level of varieties.  We now lift
this to the $\Ainf$ bar complex, recovering
Theorem~\hyperref[thm:ds-koszul-intertwine]{ds-koszul-intertwine}
from Volume~I as a consequence of Hamiltonian reduction of fiber
products.

\begin{proposition}[DS reduction as Hamiltonian reduction of the
Steinberg; \ClaimStatusConjectured]
\label{prop:ds-hamiltonian-steinberg}
Let $\cA$ be a chiral algebra with Steinberg object
$\fStein_b = \Lloop_b \times_{\Mvac} \Lloop_b$,
and let $(\mathfrak{n}, \chi)$ be DS reduction data with
$\chi \colon \mathfrak{n} \to \C$ a non-degenerate character.
The DS-reduced Steinberg is
\[
\fStein_b^{\mathrm{DS}}
  \;=\; \fStein_b /\!\!/_\chi N
  \;=\; (\Lloop_b /\!\!/_\chi N)
        \,\times_{\Mvac / N}\,
        (\Lloop_b /\!\!/_\chi N).
\]
This commutes with the convolution product, so that
\[
H_*^{\mathrm{BM}}(\fStein_b^{\mathrm{DS}})
  \;\cong\;
H_*^{\mathrm{BM}}(\fStein_b) /\!\!/_\chi N,
\]
which is the bar-cobar algebra for the DS-reduced chiral algebra
$\mathrm{DS}_\chi(\cA)$.
\end{proposition}

\begin{proof}
Two ingredients:
\begin{enumerate}[label=(\alph*)]
\item \emph{Freeness.}  The group $N$ acts freely on the regular
  locus of $\Lloop_b$: the regularity condition in
  Drinfeld--Sokolov reduction is precisely the condition that the
  stabilizer of the $\mathfrak{n}$-action on the loop space is
  trivial.  Concretely, a regular element $\varphi \in \Lloop_b$
  has $\operatorname{Stab}_N(\varphi) = \{e\}$.

\item \emph{Free actions commute with fiber products.}  When $N$
  acts freely on $\Lloop_b$, the quotient map
  $\Lloop_b \to \Lloop_b / N$ is a principal $N$-bundle.  The
  fiber product of principal bundles satisfies
  \[
  (\Lloop_b \times_{\Mvac} \Lloop_b) / N
    \;\cong\;
  (\Lloop_b / N) \times_{\Mvac / N} (\Lloop_b / N),
  \]
  and the same holds for the Hamiltonian reduction
  $/\!\!/_\chi$ (which is the quotient of the $\chi$-level set).
\end{enumerate}
The convolution product on $\fStein_b$ is defined by the
correspondence structure of the fiber product; since reduction
preserves the fiber product, it preserves convolution.  Passing to
Borel--Moore homology gives the claimed isomorphism.
\end{proof}

\begin{corollary}[$\Ainf$ DS-bar intertwining; \ClaimStatusConjectured]
\label{cor:ainf-ds-bar}
Drinfeld--Sokolov reduction preserves the $(-1)$-shifted symplectic
structure on~$\fStein_b$.  Therefore, at the $\Ainf$ level:
\[
\barB\bigl(\mathrm{DS}_\chi(\cA)\bigr)
  \;\simeq\;
\barB(\cA) /\!\!/_\chi N.
\]
In particular, the $\Ainf$ bar complex of
$\mathcal{W}_k(\mathfrak{g}) = \mathrm{DS}_\chi(V_k(\mathfrak{g}))$
equals the Hamiltonian reduction of the $\Ainf$ bar complex of
$V_k(\mathfrak{g})$.  This recovers
Theorem~\hyperref[thm:ds-koszul-intertwine]{ds-koszul-intertwine}
from Volume~I from the Steinberg perspective.
\end{corollary}

\begin{proof}
Hamiltonian reduction of a $(-1)$-shifted symplectic derived stack by
a free group action yields a $(-1)$-shifted symplectic stack (this is
the AKSZ principle: reduction preserves the shifted symplectic
form~\cite{PTVV}).  Since the $\Ainf$ structure on
$\barB(\cA)$ is encoded by the $(-1)$-shifted symplectic structure
on~$\fStein_b$
(Proposition~\ref{prop:bar-cobar-steinberg}), and since the
Hamiltonian reduction by $N$ preserves this structure, the reduced
bar complex $\barB(\cA) /\!\!/_\chi N$ carries a canonical $\Ainf$
structure.  By
Proposition~\ref{prop:ds-hamiltonian-steinberg}, the reduced
Steinberg is the Steinberg for $\mathrm{DS}_\chi(\cA)$, so
this $\Ainf$ structure is $\barB(\mathrm{DS}_\chi(\cA))$.
\end{proof}

\begin{remark}[Why DS commutes with bar-cobar at the $\Ainf$ level]
\label{rem:ds-bar-cobar-ainf}
The question ``does DS commute with bar-cobar at the $\Ainf$ level?''
receives an affirmative answer from the Steinberg perspective: it is a
consequence of the basic fact that \emph{Hamiltonian reduction commutes
with fiber products} when the group action is free.  The freeness is
not an accident but the defining regularity condition of
Drinfeld--Sokolov reduction.  From this viewpoint,
Theorem~\hyperref[thm:ds-koszul-intertwine]{ds-koszul-intertwine}
is not an algebraic miracle but a geometric tautology: the Steinberg
variety for the reduced algebra \emph{is} the reduced Steinberg
variety.
\end{remark}

\begin{remark}[Quantization parameters from spectral data]
\label{rem:quantization-spectral}
The quantization parameters of $\mathcal{W}_k(\mathfrak{g})$ are
encoded entirely in the spectral curve.  The level $k$ determines the
characteristic polynomial of the Higgs field $\varphi$:
\[
\det(\xi \cdot \mathrm{id} - \varphi)
  = \xi^N + \sum_{i=2}^{N} (-1)^i \, u_i(z) \, \xi^{N-i},
  \qquad u_i \in H^0(X, K_X^{\otimes i}),
\]
whose zero locus in $T^*X$ is the spectral curve $\Sigma_b$ for
$b = (u_2, \ldots, u_N) \in \mathcal{B}$.  Different levels give different
Hitchin base points $b$, hence different spectral curves, hence
different Steinberg varieties $\fStein_{\mathrm{Hit}}$, hence different
$\mathcal{W}$-algebras --- all from a single geometric construction.

The critical level $k = -h^\vee$ corresponds to the discriminant locus
$\mathcal{B}^{\mathrm{sing}}$ where $\Sigma_b$ acquires nodes: the
Steinberg variety degenerates, reflecting the Feigin--Frenkel center of
$\widehat{\mathfrak{g}}_{-h^\vee}$.  The self-dual point
$k + h^\vee \mapsto -(k + h^\vee)$ of the Feigin--Frenkel duality
$k \leftrightarrow -k - 2h^\vee$ acts on spectral curves by the
involution $\xi \mapsto -\xi$ of the cotangent fiber.
\end{remark}

\subsection{Quantization and loop corrections}

\subsubsection{Classical vs.\ quantum chiral algebras}

\begin{definition}[Quantum correction parameter]
In the reduction from 4d to 2d, there is a natural parameter:
\[\hbar = \epsilon_1\]

This is the $\Omega$-background parameter, which becomes Planck's constant in
the reduced theory.

Classical limit: $\hbar \to 0$ (or $\epsilon_1 \to 0$)

Quantum theory: $\hbar$ finite
\end{definition}

\begin{theorem}[Genus-graded bar complex; \ClaimStatusProvedHere]
\label{thm:genus-graded-bar}
The full bar complex has a genus-graded deformation
\[\bar{B}^{\mathrm{ch}}_\hbar(\mathcal{A}) = \bar{B}^{\mathrm{ch}}(\mathcal{A})[[\hbar]]\]
with differential
\[d_\hbar = d_0 + \hbar d_1 + \hbar^2 d_2 + \cdots\]
where $d_g$ (the genus-$g$ component of $\Dg{\bullet}$, Convention~\ref*{conv:higher-genus-differentials}) is the genus-$g$ contribution from the modular operad decomposition.
The nilpotency $(d_\hbar)^2 = 0$ holds as a formal power series identity.
\end{theorem}

\begin{proof}
By Theorem~\ref*{thm:genus-induction-strict}, the full bar complex
$\bar{B}^{\mathrm{full}}(\mathcal{A}) = \bigoplus_g \bar{B}^{(g)}(\mathcal{A})$
carries a differential whose genus-$g$ component $d_g$ (in the sense of Convention~\ref*{conv:higher-genus-differentials}) is assembled from
the modular operad structure.  Setting $\hbar^g$ as the genus bookkeeping
parameter, the relation $(d_\hbar)^2 = 0$ is the genus-by-genus
identity $\sum_{g_1 + g_2 = g} d_{g_1} \circ d_{g_2} = 0$
for each $g$, which is proved in Theorem~\ref*{thm:genus-induction-strict}.
\end{proof}

\begin{remark}[Physics interpretation]
The nilpotency $(d_\hbar)^2 = 0$ is the Maurer--Cartan equation $\ell_1(\Theta) + \tfrac{1}{2}\ell_2(\Theta, \Theta) = 0$ of Theorem~\ref*{thm:explicit-theta}, expanded genus-by-genus.  In the Costello--Li framework, $\hbar = \epsilon_1$ and $d_k$ is the $k$-loop correction; this identification requires gauge-theoretic input.
\end{remark}

\begin{example}[Two classical limits of Virasoro]
The Lie-algebraic limit $c = 0$ gives the Witt algebra (no central extension, $\Theta_\cA = 0$, uncurved bar complex).  The semiclassical limit $c \to \infty$ ($\epsilon_2 \to 0$ in AGT, $b^2 = \epsilon_1/\epsilon_2 \to \infty$) gives Poisson brackets on the Hitchin base ($\kappa \to \infty$, strongly curved $A_\infty$ structure).  These are distinct limits: $c = 0$ is uncurved, $c \to \infty$ is maximally curved.
\end{example}

%================================================================
% SECTION: W-ALGEBRAS IN BOTH FRAMEWORKS
%================================================================

\subsection{$\mathcal{W}$-algebras in both frameworks}
\label{sec:w-algebras-unifying}

\subsubsection{$\mathcal{W}$-algebras from 2D CFT perspective}

\begin{definition}[\texorpdfstring{$\mathcal{W}$}{W}-algebra (CFT definition)]\label{def:w-algebra-cft}
Following Zamolodchikov \cite{Zamolodchikov} and Fateev--Lukyanov \cite{FL88}, a
\emph{W-algebra} $\mathcal{W}$ is a vertex operator algebra containing:
\begin{enumerate}
\item Virasoro element $L$ (conformal weight 2)
\item Additional generators $W^{(s)}$ of conformal weights $s > 2$
\item Relations ensuring associativity of OPE
\end{enumerate}
\end{definition}

Standard examples: $\mathcal{W}_3$ has generators $L, W$ with $\text{wt}(L) = 2$, $\text{wt}(W) = 3$; $\mathcal{W}_N$ has generators of weights $2, 3, \ldots, N$; and $\mathcal{W}_{1+\infty}$ has infinitely many generators.

\subsubsection{$\mathcal{W}$-algebras from gauge theory perspective}

\begin{theorem}[Arakawa--Creutzig--Linshaw 2019; \ClaimStatusProvedElsewhere]\label{thm:w-from-gauge}
\cite{ACL19} Let $G$ be a simple Lie group and $\rho: G \to GL(V)$ a
representation. Consider the symplectic quotient (Higgs branch):
\[\mathcal{M}_H = \mu^{-1}(0) / G\]
where $\mu: T^*V \to \mathfrak{g}^*$ is the moment map.

Then:
\begin{enumerate}
\item The Higgs branch $\mathcal{M}_H$ carries a holomorphic symplectic structure

\item Quantization of functions $\mathcal{O}(\mathcal{M}_H)$ produces a vertex
      algebra $V_H$

\item $V_H$ contains a W-algebra $\mathcal{W}_k(\mathfrak{g})$ at level $k$
      determined by the gauge coupling

\item This matches the W-algebra from coset construction:
      \[\mathcal{W}_k(\mathfrak{g}) = \text{Com}(\mathfrak{g}_k, V_\rho)\]
\end{enumerate}
\end{theorem}


\subsubsection{Bar-cobar duality for $\mathcal{W}$-algebras}

\begin{conjecture}[\texorpdfstring{$\mathcal{W}$}{W}-algebra bar-cobar duality; \ClaimStatusConjectured]\label{conj:w-algebra-bar-cobar}
Let $\mathcal{W}_k(\mathfrak{g})$ be a W-algebra (from either construction). Then:
\begin{enumerate}
\item \textup{[\ClaimStatusProvedHere]}
The chiral envelope $\mathcal{A}_{\mathcal{W}}$ admits geometric bar
      construction:
      \[\bar{B}^{\text{geom}}(\mathcal{A}_{\mathcal{W}}) = \bigoplus_{n \geq 0}
      \Gamma\left(\overline{C}_{n+1}(X), \mathcal{A}_{\mathcal{W}}^{\boxtimes(n+1)}
      \otimes \Omega^n_{\log}\right)\]

\item At generic level $k$ \textup{[\ClaimStatusProvedHere]};
      at admissible levels \textup{[\ClaimStatusConjectured]}:
      $\mathcal{W}_k(\mathfrak{g})$ is \emph{Koszul} and has a chiral
      Koszul dual coalgebra $\mathcal{A}_{\mathcal{W}}^!$

\item Conditional on~\textup{(2)}: the bar and cobar complexes are quasi-inverse:
      \[\Omega(\bar{B}(\mathcal{A}_{\mathcal{W}})) \simeq \mathcal{A}_{\mathcal{W}}\]
      At generic $k$ this is \textup{[\ClaimStatusProvedHere]};
      at admissible levels \textup{[\ClaimStatusConjectured]}

\item \textup{[\ClaimStatusProvedHere]}
All structures (Virasoro, W-currents, OPE) have geometric realization via
      configuration spaces
\end{enumerate}
\end{conjecture}

\begin{proof}[Proof of items \textup{(1)} and \textup{(4)}]
Item~(1): the geometric bar construction applies to any augmented chiral algebra
(Definition~\ref*{def:geometric-bar}); the bar differential satisfies
$d^2 = 0$ by Theorem~\ref*{thm:genus-induction-strict}.  For W-algebras,
the chiral envelope exists by the Beilinson--Drinfeld functor
(Theorem~3.7.11 of \cite{BD04}).

Item~(4): OPE data is encoded by residues at collision divisors of the
FM compactification (Definition~\ref*{def:bar-differential-complete}).
For W-algebra currents $W^{(s)}(z)$, the OPE poles of orders
$\leq 2s$ correspond to the stratification of $\overline{C}_2(X)$
along the diagonal, with each pole order contributing to a specific
component of $d_{\mathrm{bar}}$.

Item~(3) follows from
Theorem~\ref*{thm:bar-cobar-isomorphism-main} once the Koszul property~(2)
is established.
\end{proof}

\begin{proof}[Proof of items \textup{(2)} and \textup{(3)} at generic~$k$]
Let $k$ be generic (away from admissible and critical levels).

\emph{Step~1: Koszul property via weight spectral sequence.}
At generic~$k$, the W-algebra $\mathcal{W}_k(\mathfrak{g})$ is simple
as a vertex algebra (Arakawa \cite{Ara07}, Theorem~4.1).  Simplicity
implies that Verma modules are irreducible
(Proposition~\ref*{prop:generic-irreducibility}; the Shapovalov form
is non-degenerate at each weight), so
$\mathrm{Ext}^q_{\mathcal{O}_k}(V(h_1), V(h_2)) = 0$ for $q > 0$
and all pairs of distinct Verma modules.

The weight filtration on the bar complex induces a spectral sequence
$E_2^{p,q} \Rightarrow \ChirHoch^{p+q}(\mathcal{W}_k(\mathfrak{g}))$.
The Ext-vanishing forces $E_2^{p,q} = 0$ for $q > 0$, so the spectral
sequence collapses at~$E_2$.  This $E_2$-degeneration is equivalent to
the Koszul property: the bar resolution is linear
(Adams-graded components are acyclic in positive internal degree).

\emph{Step~2: Identification of the Koszul dual.}
By Theorem~\ref*{thm:w-algebra-hochschild} (\ClaimStatusProvedHere),
the bar complex cohomology at generic~$k$ is:
\[
\ChirHoch^*(\mathcal{W}_k(\mathfrak{g}))
\;=\; \mathbb{C}[\Theta_1, \ldots, \Theta_r],
\qquad \deg \Theta_i = h_i = m_i + 1
\]
where $m_1, \ldots, m_r$ are the exponents of~$\mathfrak{g}$.
The Koszul dual coalgebra $\mathcal{A}_{\mathcal{W}}^!$ is determined by
this cohomology: it is the cofree coalgebra cogenerated by the classes
$\Theta_i^*$ dual to the $\Theta_i$.

\emph{Step~3: Bar-cobar quasi-isomorphism.}
With the Koszul property established,
Theorem~\ref*{thm:bar-cobar-isomorphism-main} gives the
quasi-isomorphism
$\Omega(\bar{B}(\mathcal{A}_{\mathcal{W}})) \simeq
\mathcal{A}_{\mathcal{W}}$.
\end{proof}


\begin{example}[Explicit: \texorpdfstring{$\mathcal{W}_3$}{W_3} at \texorpdfstring{$c = -2$}{c = -2}]\label{ex:w3-c-minus-2-explicit}
For the $\mathcal{W}_3$ algebra at central charge $c = -2$:

\emph{Generators.}
\begin{align*}
L &= \text{energy-momentum tensor, conformal weight 2} \\
W &= \text{W-current, conformal weight 3}
\end{align*}

\emph{OPE.}
\begin{align*}
L(z)L(w) &\sim \frac{c/2}{(z-w)^4} + \frac{2L(w)}{(z-w)^2} + \frac{\partial L(w)}{z-w} \quad (c = -2) \\
L(z)W(w) &\sim \frac{3W(w)}{(z-w)^2} + \frac{\partial W(w)}{z-w} \\
W(z)W(w) &\sim \frac{c_{33}}{(z-w)^6} + \frac{2L(w)}{(z-w)^4} + \cdots
\end{align*}
where $c_{33} = c/3 = -2/3$ by the standard Zamolodchikov normalization.

\emph{Chiral algebra presentation.}
\[\mathcal{A}_{\mathcal{W}_3} = \text{Free}_{\mathcal{D}}(\mathbb{C}L \oplus \mathbb{C}W)
/ \text{(W-algebra relations)}\]

\emph{Bar complex degree~2.}
\[\bar{B}^2(\mathcal{A}_{\mathcal{W}_3}) = \Gamma\left(\overline{C}_3(X),
\mathcal{A}_{\mathcal{W}_3}^{\boxtimes 3} \otimes \Omega^\bullet\right)\]

Elements are represented by integrals:
\[\int_{\overline{C}_3(X)} f(z_1, z_2, z_3) \wedge d\log(z_1-z_2) \wedge d\log(z_2-z_3)\]
where $f$ is a section of $\mathcal{A}_{\mathcal{W}_3}^{\boxtimes 3}$.

\emph{Differential action.}
\begin{align*}
d_{\text{mult}}(f) &= \text{Res}_{z_1 = z_2}[f] + \text{Res}_{z_2 = z_3}[f] \\
&\quad + \text{Res}_{z_1 = z_3}[f] \quad \text{(collisions)} \\
d_{\text{internal}}(f) &= d_{\mathcal{W}}(f) \quad \text{(internal differential)} \\
d_{\text{extend}}(f) &= \text{extension across boundary divisors}
\end{align*}

\emph{Koszul dual coalgebra.} At $c = -2$, we conjecture (building on Arakawa's rationality result~\cite{Ara07}) that $\mathcal{W}_3$ is Koszul, with coalgebra dual $\mathcal{C}_{\mathcal{W}_3}$ given by:
\[\mathcal{C}_{\mathcal{W}_3} \subset \text{Cofree}_{\mathcal{D}}(s^{-1}\mathbb{C}L^* \oplus s^{-1}\mathbb{C}W^*)\]
the sub-coalgebra cogenerated by the orthogonal complement of the relations, where $s^{-1}$ denotes desuspension.
\end{example}

%================================================================
% SECTION: MATHEMATICAL BRIDGES
%================================================================

\subsection{AGT correspondence and bar-cobar duality}
\label{sec:mathematical-bridges}



\subsubsection{AGT correspondence via bar-cobar}

\begin{theorem}[AGT 2D side: bar complex = semi-infinite complex; \ClaimStatusProvedHere]
\label{thm:agt-2d-bar}
\index{AGT correspondence!2D side}
Let $\mathcal{W}_k(\fg)$ be a W-algebra at generic level~$k$.  Then $H^\bullet(\barB^{\mathrm{ch}}(\mathcal{W}_k(\fg))) \cong H^{\frac{\infty}{2}+\bullet}(\mathcal{W}_k(\fg))$ (Theorem~\ref*{thm:bar-semi-infinite-w}), the anomaly duality $\kappa(\mathcal{W}_k(\fg)) + \kappa(\mathcal{W}_{k'}(\fg^\vee)) = 0$ holds for Feigin--Frenkel dual pairs (Corollary~\ref*{cor:anomaly-duality-w}), and chiral homology is computed by configuration space integrals on FM compactifications (Theorem~\ref*{thm:bar-NAP-homology}).
\end{theorem}

\begin{conjecture}[AGT 4D--2D bridge via bar-cobar; \ClaimStatusConjectured]\label{conj:agt-bar-cobar}
\index{AGT correspondence}
The Alday--Gaiotto--Tachikawa correspondence \cite{AGT09} admits a
bar-cobar formulation: the Costello--Li holomorphic twist of 4D
$\mathcal{N}=2$ gauge theory with gauge group~$G$ produces a
factorization algebra $\mathcal{F}_G$ on a curve~$\Sigma$, and
this factorization algebra is equivalent to $\mathcal{W}_k(\fg)$:

\begin{center}
\begin{tikzcd}
\text{4D } \mathcal{N}=2 \text{ gauge partition function} \ar[d, "\text{CL twist}"] \ar[r, "\text{AGT}"]
& \text{2D Liouville/Toda CFT correlation} \ar[d, "\text{chiral envelope}"] \\
\text{2D factorization algebra } \mathcal{F}_G \ar[r, "\simeq"]
& \text{W-algebra } \mathcal{W}_k(\mathfrak{g})
\end{tikzcd}
\end{center}

Moreover:
\begin{enumerate}
\item The 4D instanton partition function = W-algebra conformal blocks
\item The 4D Coulomb branch parameters = W-algebra momenta
\item The BV complex of the 4D theory computes the same homology as
  the bar complex of the 2D chiral algebra:
      \[H_\bullet^{\text{ch}}(X, \mathcal{W}_k(\mathfrak{g})) = H_\bullet^{\text{BV}}(\mathcal{F}_G)\]
\end{enumerate}
\end{conjecture}

\begin{remark}[Scope]
The 2D side is proved (Theorem~\ref{thm:agt-2d-bar}); the 4D--2D bridge is proved in specific cases \cite{SV13,MO19}.  MC4 and MC5 are both proved in Volume~I (Chapter~\ref*{chap:concordance}); the remaining obstruction is the 4D--2D bridge for general gauge groups.
\end{remark}

%----------------------------------------------------------------
\subsubsection{The instanton Steinberg and the 4D/2D bridge}
\label{sec:instanton-steinberg}
\index{Steinberg variety!instanton}
\index{AGT correspondence!Steinberg presentation}
%----------------------------------------------------------------

\providecommand{\Minst}{\mathcal{M}_{\mathrm{inst}}}
\providecommand{\Sinst}{\mathfrak{S}_{\mathrm{inst}}}

The Chriss--Ginzburg philosophy --- that representation-theoretic
objects arise as Borel--Moore convolution algebras on Steinberg-type
correspondences --- applies directly to the AGT correspondence, yielding
a structural explanation rather than a case-by-case verification.

\begin{construction}[Instanton Steinberg correspondence]
\label{constr:instanton-steinberg}
\index{instanton moduli!Steinberg correspondence}
Let $G$ be a reductive group and fix the framing data at infinity on
$\mathbb{C}^2$.  The moduli space $\Minst(G,n)$ of framed $G$-instantons
of charge~$n$ on $\mathbb{C}^2$ carries a natural Hecke correspondence
via the framed moduli $\Minst^0$:
\begin{equation}\label{eq:instanton-steinberg}
  \Sinst \;=\; \Minst(G,n) \;\times_{\Minst^0}\; \Minst(G,n).
\end{equation}
The Borel--Moore convolution algebra
\[
  H_*^{\mathrm{BM}}(\Sinst) \;\cong\; \mathcal{W}_k(\fg)
\]
recovers the W-algebra, where the level~$k$ is determined by the
$\Omega$-background parameters $(\varepsilon_1, \varepsilon_2)$ via
$k + h^\vee = -\varepsilon_1/\varepsilon_2$.  This is the
Schiffmann--Vasserot \cite{SV13} and Maulik--Okounkov \cite{MO19}
theorem: the W-algebra is the \emph{instanton Steinberg algebra},
precisely as the affine Hecke algebra is the Steinberg algebra of
$\widetilde{\mathcal{N}} \times_{\mathcal{N}} \widetilde{\mathcal{N}}$.
\end{construction}

\begin{proposition}[AGT as Chriss--Ginzburg; \ClaimStatusProvedElsewhere{} \cite{SV13,MO19}]
\label{prop:agt-chriss-ginzburg}
\index{AGT correspondence!Chriss--Ginzburg interpretation}
The AGT correspondence
$Z_{\mathrm{Nekrasov}} = \langle \text{conformal block} \rangle$
is the instanton analog of the Chriss--Ginzburg character formula.
Specifically:
\begin{enumerate}
\item The Nekrasov partition function is the pushforward of the
  fundamental class $[\Minst(G,n)]$ along the projection
  $\Minst(G,n) \to \mathrm{pt}$, computed by equivariant localization.
\item The conformal block is the matrix element of the convolution
  algebra $H_*^{\mathrm{BM}}(\Sinst) \cong \mathcal{W}_k(\fg)$
  between Whittaker states.
\item Their equality is the direct analog of the classical statement
  that the character of an $H_W$-module equals the pushforward of the
  fundamental class of the corresponding Springer fiber.
\end{enumerate}
Thus AGT is not a mysterious coincidence between 4D and 2D physics but
the Steinberg presentation theorem applied to instanton moduli.
\end{proposition}

\begin{proof}
The Schiffmann--Vasserot/Maulik--Okounkov identification
$H_*^{\mathrm{BM}}(\Sinst) \cong \mathcal{W}_k(\fg)$ provides the
algebra isomorphism.  Under this isomorphism, the equivariant
fundamental class $[\Minst(G,n)]^T$ maps to the Whittaker vector
$|W\rangle$ in the Verma module (this is the fixed-point basis
correspondence of \cite{MO19}).  The Nekrasov partition function
\[
  Z_{\mathrm{Nek}}(\vec{a}, \varepsilon_1, \varepsilon_2, \vec{m}, q)
  \;=\; \sum_n q^n \int_{[\Minst(G,n)]^T} 1
\]
is the equivariant pushforward, while the conformal block
$\langle W | \Phi_1(z_1) \cdots \Phi_m(z_m) | W \rangle$ is the
matrix element of vertex operators in the convolution algebra.  The
equality of these two expressions is exactly the content of the
Chriss--Ginzburg localization theorem applied to the instanton Steinberg
variety: the character of a representation (= matrix element) equals
the pushforward from the fiber (= localized integral).
\end{proof}

\begin{remark}[Bar-cobar as 4D/2D dictionary]
\label{rmk:bar-cobar-4d-2d}
\index{bar-cobar adjunction!4D/2D interpretation}
The bar-cobar adjunction acquires a transparent 4D/2D meaning through
the Steinberg lens:
\begin{itemize}
\item \textbf{Bar = open-string/instanton partition function.}  Counting
  instantons with boundary conditions is counting bar elements on
  Fulton--MacPherson configurations: the bar differential $d_B$ from OPE
  residues on $\mathrm{FM}_k(\mathbb{C})$ encodes the instanton
  interaction vertices, and the bar degree~$k$ is the instanton
  charge.
\item \textbf{Cobar = 2D CFT recovery.}  Reconstructing the W-algebra
  from its bar coalgebra data $\Omega(\barB(\mathcal{W}_k(\fg)))
  \simeq \mathcal{W}_k(\fg)$ is reconstructing the 2D conformal field
  theory from its instanton sector.  The cobar construction is the
  inverse Steinberg presentation: from convolution data, recover the
  algebra.
\end{itemize}
The bar-cobar adjunction of Volume~I, Theorem~A, specialized to the
instanton setting, \emph{is} the 4D/2D correspondence.
\end{remark}

\begin{remark}[Costello--Gaiotto and the HT Steinberg]
\label{rmk:costello-gaiotto-steinberg}
\index{Costello--Gaiotto!instanton Steinberg}
The Costello--Gaiotto program \cite{CG18} provides the bridge between
the physical 4D $\mathcal{N}=2$ theory and the algebraic Steinberg
picture.  The holomorphic twist of the 4D theory on
$\mathbb{C}_z \times \mathbb{C}_w$ produces, upon compactification on
$\mathbb{C}_w$ with $\Omega$-background, a holomorphic-topological
theory on $\mathbb{C}_z \times \mathbb{R}_t$ whose bar complex is
precisely the instanton Steinberg variety $\Sinst$.  The Steinberg
presentation theorem (that $H_*^{\mathrm{BM}}(\Sinst)$ is the
W-algebra) then gives the AGT correspondence as a \emph{special case}
of the algebraic Steinberg presentation --- the same mechanism by which
the bar complex presents the Swiss-cheese algebra in Part~I.
\end{remark}



%================================================================
% SECTION: OPEN QUESTIONS
%================================================================

\subsection{Open questions and future directions}
\label{sec:open-questions-holo-topo}

\begin{theorem}[\texorpdfstring{$E_n$}{E_n} Koszul duality on \texorpdfstring{$n$}{n}-manifolds;
\ClaimStatusProvedElsewhere]\label{thm:en-koszul}
\index{En algebra@$E_n$ algebra!Koszul duality}
\index{Koszul duality!higher-dimensional}
For an $E_n$-algebra $\mathcal{A}$ on a smooth oriented $n$-manifold $M$,
there is a bar-cobar adjunction
\begin{equation}\label{eq:en-bar-cobar}
\barB_{E_n}: E_n\text{-}\mathrm{Alg}(M)
\rightleftarrows
E_n\text{-}\mathrm{CoAlg}(M) :\Omega_{E_n}
\end{equation}
realized by configuration space integrals on the Axelrod--Singer
compactification $\overline{C}_k(M)$ of $k$-point configurations in $M$.
The propagator is the Green kernel of the Hodge Laplacian on $M$, and
the Arnold relations are replaced by the Totaro relations
\cite{Totaro96} for $H^*(\overline{C}_k(\mathbb{R}^n))$.

For $n = 1$, this recovers the classical (non-chiral) bar-cobar
adjunction for $A_\infty$-algebras.  For $n = 2$ on a Riemann surface,
it recovers the chiral bar-cobar adjunction of
Theorem~\textup{\ref*{thm:bar-cobar-isomorphism-main}}.
The Poincar\'e--Koszul duality of Ayala--Francis~\cite{AF15}
\textup{(}Theorem~7.8\textup{)} provides the abstract framework;
our propagator formulas should supply explicit chain-level
representatives.

\emph{Scope.}  For $n \geq 3$, the configuration space integrals involve
the Axelrod--Singer propagator (which is not holomorphic), and the
combinatorics of boundary strata are governed by graph complexes
(Kontsevich \cite{Kon94}).  The key missing ingredient is an
explicit form-level presentation of the $E_n$ bar differential in terms
of iterated integrals over $\overline{C}_k(M)$.  The $n = 3$ case
with $M = S^3$ connects to the Chern--Simons perturbative expansion and
Vassiliev invariants (see Remark~\textup{\ref{rem:vassiliev}} below).

\emph{Resolution.}  All four components are now
established: see Theorem~\textup{\ref*{thm:en-koszul-duality}}, which
assembles the proofs of Totaro, Lambrechts--Voli\'c, Knudsen, and
our $n = 2$ recovery.
\end{theorem}


\begin{remark}[Vassiliev invariants from Feynman transform]
\label{rem:vassiliev}
\index{Vassiliev invariants!from bar complex}
The Prism Principle $\barB^{\mathrm{full}} \simeq \mathrm{FT}_{\mathrm{Mod}}$ (Theorem~\ref*{thm:prism-higher-genus}(ii)) and Getzler--Kapranov \cite{GeK98} suggest (\ClaimStatusConjectured) that the bar complex of the Chern--Simons factorization algebra on $S^3$ produces universal Vassiliev invariants, recovering the Kontsevich integral via specialization to the real locus.  The main gap is the passage $\overline{C}_n(X) \to \overline{C}_n(S^1)$.
\end{remark}

%================================================================
% SECTION: BV-HT PHYSICS
% (from bv_ht_physics.tex)
%================================================================

\section{Holomorphic-topological field theories and BV-BRST structure}
\label{sec:bv-ht-physics}

\subsection{Gaiotto's framework: from 4d to 2d}

\begin{definition}[Holomorphic-topological (HT) theory]
A holomorphic-topological field theory on a complex surface $\Sigma$ is holomorphic in one complex direction and topological in the other.  Fields are sections of $\mathcal{O}$-modules satisfying $\bar{\partial}_{\bar{z}} \phi = 0$ (holomorphic) and $Q_{\text{top}} \phi = 0$ (topologically BRST-closed).
\end{definition}

\begin{theorem}[HT theory from 4d \texorpdfstring{$\mathcal{N}=4$}{N=4} SYM; \ClaimStatusProvedElsewhere]
\label{thm:ht-from-n4-sym}
Starting with 4d $\mathcal{N}=4$ super Yang--Mills with gauge group $G$, one applies the holomorphic twist (Costello~\cite{costello-gaiotto}), localizes to $\bar{\partial}$-connections on a Riemann surface $\Sigma$, and obtains holomorphic BF theory (the dimensional reduction of holomorphic Chern--Simons to $\Sigma$).

The action is:
\[S_{\text{BF}} = \int_{\Sigma} \operatorname{Tr}\!\left(B \wedge \bar{\partial} \mathbf{A} + B \wedge \mathbf{A} \wedge \mathbf{A}\right)\]

where $\mathbf{A} = c + A + \cdots$ denotes the full BV field with $c \in \Omega^{0,0}(\Sigma, \mathfrak{g})$ (ghost) and $A \in \Omega^{0,1}(\Sigma, \mathfrak{g})$ (connection), and $B \in \Omega^{1,0}(\Sigma, \mathfrak{g})$ is the adjoint-valued Lagrange multiplier.
On a Riemann surface, $\Omega^{0,2}(\Sigma) = 0$, so the classical fields
$A \in \Omega^{0,1}$ satisfy $\bar{\partial}A = 0$ and $A \wedge A = 0$ identically.  The non-trivial
content of the action arises from the BV extension: the ghost component $c \in \Omega^{0,0}$ has
$\bar{\partial}c \in \Omega^{0,1}$, giving $B \wedge \bar{\partial}c \in \Omega^{1,1}$ (the volume form).
The full holomorphic Chern--Simons action $\int \Omega \wedge \operatorname{CS}(\mathcal{A})$ lives on a
complex $3$-fold with $(3,0)$-form $\Omega$; on $\Sigma$ only the holomorphic BF action above survives the dimensional reduction.
\end{theorem}

\begin{proof}[From Costello--Gaiotto]
\emph{Step~1: The Holomorphic Twist.}

Start with $\mathcal{N}=4$ SYM on $\mathbb{R}^4 \cong \mathbb{C}^2$. The field
content includes the gauge field $A_\mu$, scalars $\Phi^I$ in the adjoint ($I=1,\ldots,6$), and fermions $\psi, \bar{\psi}$.  The twist redefines the Lorentz group action by mixing with R-symmetry:
\[\text{Spin}(4) \times \text{Spin}(6)_R \to \text{Spin}(4)_{\text{new}}\]

After twisting, some bosons become fermions (ghosts), some fermions become bosons (matter), and a nilpotent supercharge $Q$ becomes scalar.

\emph{Step~2: Localization.}

The twisted action has $Q^2 = 0$ and:
\[S_{\text{twisted}} = \{Q, \Lambda\} + S_0\]

By localization, path integral reduces to:
\[Z = \int_{\text{Q-fixed locus}} \mathcal{O}(fields)\]

The Q-fixed locus consists of holomorphic data:
\[\bar{\partial}_A = 0, \qquad \bar{\partial}_A \Phi = 0\]

These are the equations for a holomorphic bundle with holomorphic Higgs field (the Hitchin equations in the holomorphic sector).

\emph{Step~3: Volume Form.}

The holomorphic volume form $\Omega$ arises from the topological twist. On
$\mathbb{C}^2$ with holomorphic coordinates $(z,w)$:
\[\Omega = dz \wedge dw\]

On the deformed conifold $Y_\mu = \{u_1 w_2 - u_2 w_1 = \mu\} \subset \mathbb{C}^4$ (a Calabi--Yau 3-fold), $\Omega$ is the unique holomorphic $(3,0)$-form obtained via the Poincar\'{e} residue:
\[\Omega = \operatorname{Res}_{Y_\mu}\frac{du_1 \wedge du_2 \wedge dw_1 \wedge dw_2}{u_1 w_2 - u_2 w_1 - \mu}\]
\end{proof}

\subsection{Boundary conditions and chiral algebras}

\begin{theorem}[Boundary chiral algebra; \ClaimStatusProvedElsewhere{} \cite{CDG20,CG17}]
\label{thm:boundary-chiral-algebra-bv}
A boundary condition for holomorphic Chern--Simons theory supports a chiral
algebra $\mathcal{A}_{\text{bdy}}$ whose generators are local operators at the boundary, whose OPE comes from bulk-to-boundary correlation functions, and whose central charge is determined by the level of HCS theory.
\end{theorem}

\begin{example}[Kac--Moody from HCS]
Holomorphic Chern--Simons with gauge group $G$ at level $k$ produces:
\[\mathcal{A}_{\text{bdy}} = \widehat{\mathfrak{g}}_k\]

the affine Kac--Moody algebra at level $k$.

The currents:
\[J^a(z) = \text{Tr}(T^a A(z))\]

satisfy:
\[J^a(z) J^b(w) \sim \frac{k \delta^{ab}}{(z-w)^2} +
\frac{f^{abc} J^c(w)}{z-w}\]

This OPE is computed via the holomorphic Chern--Simons path integral with
boundary insertions.
\end{example}

\subsection{The holomorphic-topological boundary condition}

\begin{definition}[HT boundary condition]
For holomorphic Chern--Simons on a surface $\Sigma$ with boundary $\partial \Sigma$,
the \emph{holomorphic-topological boundary condition} requires:
\begin{enumerate}
\item Fields extend holomorphically to a compactification $\bar{\Sigma}$
\item At the boundary divisor $D = \bar{\Sigma} \setminus \Sigma$, fields have
a simple zero: $A \in \Omega^{0,*}(\bar{\Sigma}, \mathcal{O}_{\bar{\Sigma}}(-D))$
\item The volume form $\Omega$ has compatible pole: $\Omega \in K_{\bar{\Sigma}}(kD)$
for cubic interaction $(k=3)$
\end{enumerate}
\end{definition}

\begin{conjecture}[Bar-cobar from HT boundary; \ClaimStatusHeuristic]
\label{conj:bar-cobar-ht-boundary}
The holomorphic-topological boundary condition realizes bar-cobar duality:
\begin{itemize}
\item \emph{Bar side.} Fields with prescribed asymptotics near $D$ (logarithmic
forms)
\item \emph{Cobar side.} Distributional fields on $\Sigma$ (delta functions at $D$)
\item \emph{Duality.} Perfect pairing via residue-distribution integral
\end{itemize}
\end{conjecture}

\begin{evidence}[Evidence: Geometric realization]
The key is the volume form behavior. If $\Omega$ has a pole of order $k$ along
$D$:
\[\Omega \sim \frac{dz \wedge dw}{(z-w)^k}\]

then the action term:
\[\int_{\Sigma} \Omega \cdot A^k\]

is finite if and only if $A$ vanishes to order 1 along $D$.

This pole-zero compatibility is exactly the bar-cobar relationship: the logarithmic form $\eta = d\log(z-w)$ (bar) has a logarithmic singularity whose residue is the distribution $\delta(z-w)$ (cobar).

The holomorphic-topological boundary condition enforces this duality at the
geometric level.
\end{evidence}

\begin{remark}[Scope]
This conjecture is labeled \ClaimStatusHeuristic{} because it proposes a geometric
realization of bar-cobar duality via holomorphic-topological boundary conditions,
which requires input from extended field theory (boundary conditions for HT theories)
beyond the algebraic framework of this monograph.

\emph{Homotopy template} (Convention~\ref*{conv:hms-levels}): Type~VII (physics-dictionary).
The bar and cobar complexes should arise as the two boundary conditions (logarithmic forms vs.\ distributional fields) of a holomorphic-topological theory on~$\Sigma$, with the bar-cobar pairing realized by the residue-distribution integral across the boundary divisor~$D$.
\end{remark}

\subsection{$\mathcal{W}$-algebras from Higgs branches}

\subsubsection{\texorpdfstring{4d gauge theory $\to$ 2d $\mathcal{W}$-algebra}{4d gauge theory -> 2d $\mathcal{W}$-algebra}}

\begin{conjecture}[Costello--Gaiotto AGT; \ClaimStatusConjectured]
\label{conj:costello-gaiotto-agt}
Starting with 4d $\mathcal{N}=2$ gauge theory with gauge group $G$, compactify on $\Sigma_g$, apply the topological twist, and take the infrared limit.  The result is a 2d CFT with W-algebra symmetry $\mathcal{W}(G)$.

For $G = SU(N)$, this gives the $\mathcal{W}_N$ algebra.
\end{conjecture}

\begin{evidence}[Evidence: Via bar-cobar]
\emph{Step~1: Higgs Moduli Space.}

The 4d theory has Higgs branch moduli space:
\[\mathcal{M}_{\text{Higgs}} = \text{Higgs}(\Sigma_g, G)\]

consisting of $G$-Higgs bundles on $\Sigma_g$.

\emph{Step~2: Chiral Algebra.}

The Higgs moduli space supports a chiral algebra via:
\[\mathcal{A} = \text{LocalObs}(\mathcal{M}_{\text{Higgs}})\]

Local operators on the moduli space form a factorization algebra, which extends
to a chiral algebra.

\emph{Step~3: W-Algebra Identification.}

For $G = SU(N)$ and $\Sigma_g = \mathbb{C}$, the local operators --- the stress tensor $T(z)$ from diffeomorphisms and higher spin currents $W^{(s)}(z)$ for $s = 2, \ldots, N$ --- generate the $\mathcal{W}_N$ algebra.

\emph{Step~4: Bar-Cobar Realization.}

The bar complex:
\[\bar{B}^{\text{ch}}(\mathcal{W}_N) = \Omega^*(\overline{C}_n(\Sigma_g),
\mathcal{W}_N^{\boxtimes n})\]

computes the boundary W-algebra correlators on the 2d side.  The
conjectural gauge-theory statement is that these same correlators arise
from partition functions of the 4d theory:
\[\langle W^{(s_1)}(z_1) \cdots W^{(s_n)}(z_n) \rangle =
Z_{4d}[\Sigma_g; z_1, \ldots, z_n]\]
\end{evidence}

\begin{remark}[Scope]
\ClaimStatusConjectured{} because the 4d gauge theory origin of W-algebras (Costello--Gaiotto) requires the holomorphic twist and infrared localization, which are physics inputs beyond the scope of this monograph.  The bar complex computing W-algebra conformal blocks follows from Part~I; the live task is the gauge-theoretic comparison.

\emph{Homotopy template} (Convention~\ref*{conv:hms-levels}): Type~VII (physics-dictionary).  The remaining obstacle is the gauge-theoretic comparison (physics input).
\end{remark}

\subsection{Quantum corrections and central charge}

\begin{remark}[Quantum vs.\ classical]
The 4d $\to$ 2d reduction involves two types of quantum corrections: \emph{loop corrections} from integrating out massive modes (captured by $m_3, m_4, \ldots$ in the $A_\infty$ structure), and \emph{instanton corrections} from non-perturbative effects (not visible in the bar complex, requiring full QFT).
\end{remark}

\begin{theorem}[Central charge from 4d; \ClaimStatusProvedElsewhere]
\label{thm:central-charge-from-4d}
The Sugawara central charge of the affine Kac--Moody algebra $\widehat{\mathfrak{g}}_k$ is:
\[c_{\mathrm{Sug}} = \frac{k \dim \mathfrak{g}}{k + h^\vee}\]

The W-algebra $\mathcal{W}^k(\mathfrak{g}, f)$ obtained by DS reduction has a different central charge, determined by $c_{\mathrm{Sug}}$ minus the ghost contribution (which depends on the nilpotent orbit of $f$), where $k$ is the level and $h^\vee$ the dual Coxeter number.  This matches the Feigin--Frenkel formula~\cite{Feigin-Frenkel}.
\end{theorem}

\subsection{Quantum observables and BV integration}

\subsubsection{BV path integral}

\begin{definition}[BV partition function]
The BV partition function is:
\[Z_{\text{BV}} = \int_{\mathcal{L}} [D\phi] \, e^{S[\phi]/\hbar}\]

where $\mathcal{L}$ is a Lagrangian submanifold (gauge fixing), $S[\phi]$ satisfies the quantum master equation $\Delta_{\text{BV}} e^{S/\hbar} = 0$, and the integration uses the BV measure (Berezinian).
\end{definition}

\begin{conjecture}[BV integration = bar-cobar pairing; \ClaimStatusHeuristic]
\label{conj:bv-integration-bar-cobar}
The BV path integral is realized by the bar-cobar pairing:
\[Z_{\text{BV}} = \langle \bar{B}^{\text{ch}}(\mathcal{A}),
\Omega^{\text{ch}}(\mathcal{C}) \rangle\]

Explicitly:
\[Z = \int_{\overline{C}_n(X)} \omega_{\text{bar}} \wedge \iota^* K_{\text{cobar}}\]
\end{conjecture}

\begin{evidence}
\emph{Step~1: Gauge Fixing as Lagrangian.}

Choose gauge fixing Lagrangian $\mathcal{L}$ corresponding to:
\[\mathcal{L} = \{(\phi, \phi^+) : \phi^+ = F(\phi)\}\]

for some gauge fermion $F$.

In geometric terms, this corresponds to choosing a regularization prescription
for the configuration space integrals.

\emph{Step~2: BV Measure.}

The BV measure on $\mathcal{L}$ is:
\[\mu_{\text{BV}} = \text{Ber}(\mathcal{L}) \cdot d\phi\]

Geometrically, this is the measure on configuration space:
\[\mu_{\text{geom}} = \prod_{i<j} |z_i - z_j|^2 \prod_{i=1}^n d^2z_i\]

with appropriate gauge fixing factors.

\emph{Step~3: Action and Pairing.}

The action:
\[e^{S/\hbar} = \prod_{\text{interactions}} e^{V_k(z_1, \ldots, z_k)/\hbar}\]

corresponds to the cobar element:
\[K_{\text{cobar}} = \sum_n K_n(z_1, \ldots, z_n)\]

The pairing:
\[\int_{\overline{C}_n(X)} \omega_{\text{bar}} \wedge K_{\text{cobar}}\]

is exactly the BV path integral with gauge fixing determined by the choice of
regularization.
\end{evidence}

\begin{remark}[Scope]
This conjecture is labeled \ClaimStatusHeuristic{} because the identification of the
BV path integral with the bar-cobar pairing requires a rigorous measure-theoretic
framework for configuration space integrals with BV gauge-fixing data, which goes
beyond the algebraic constructions of Part~I.

\emph{Homotopy template} (Convention~\ref*{conv:hms-levels}): Type~VII (physics-dictionary).
The BV partition function $Z_{\mathrm{BV}}$ should equal the bar-cobar pairing $\langle \bar{B}^{\mathrm{ch}}(\cA), \Omega^{\mathrm{ch}}(\mathcal{C}) \rangle$ as a configuration space integral, with gauge-fixing Lagrangian corresponding to the choice of regularization of the FM compactification.
\end{remark}

\subsubsection{From heuristic to geometry: the correspondence integral}

The heuristic of Conjecture~\ref{conj:bv-integration-bar-cobar} admits a precise
geometric upgrade, modeled on the Chriss--Ginzburg convolution formalism for
Steinberg varieties~\cite{CG97}.

\providecommand{\Steinberg}{\mathfrak{S}}
\providecommand{\BM}{\mathrm{BM}}

\begin{construction}[The correspondence integral]
\label{constr:correspondence-integral}
Let $\Steinberg_b$ denote the Steinberg object associated to the bar complex
$\barB(\cA)$ (Definition~\ref*{def:geometric-steinberg}).
Given a bar element $\beta \in \barB(\cA)$, viewed as a Borel--Moore chain on
$\Steinberg_b$, and a cobar element $\omega \in \Omega(\cA^!)$, viewed as a
cohomology class on $\Steinberg_b$, the \emph{correspondence integral} is the
Borel--Moore/cohomology pairing:
\[\langle \beta, \omega \rangle_{\Steinberg_b}
\;=\; \int_{\Steinberg_b} \beta \cdot \omega
\;\in\; H_*^{\BM}(\Steinberg_b) \otimes_{H^*(\Steinberg_b)} \mathbf{k}.\]
The BV action $S_{\mathrm{BV}}$ is the $(-1)$-shifted symplectic form
$\omega_{-1}$ on $\Steinberg_b$, and the exponential $e^{S_{\mathrm{BV}}}$
is the fundamental class $[\Steinberg_b] \in H_*^{\BM}(\Steinberg_b)$.
\end{construction}

\begin{proposition}[\ClaimStatusProvedHere]
\label{prop:bv-bar-cobar-betagamma}
For the free $\beta\gamma$ system, the BV path integral
$\int e^{S_{\mathrm{BV}}} \, D\phi$ restricted to $k$-point configurations
equals the bar-cobar pairing on
$\barB_k(\mathrm{Heis}) \otimes \Omega_k(\mathrm{Heis}^!)$.
\end{proposition}

\begin{proof}
The BV integral localizes to the critical points of $S_{\mathrm{BV}}$ on each
Fulton--MacPherson stratum $\overline{\mathrm{FM}}_k(\mathbb{C})$.  These
critical points are exactly the bar elements --- the solutions to the
Maurer--Cartan equation on
$\mathrm{FM}_k(\mathbb{C}) \times \mathrm{Conf}_k(\mathbb{R})$ --- since the
$\beta\gamma$ action is quadratic and the OPE is first-order.

The fluctuation determinant around each critical point is the ratio of
regularized Gaussian integrals on the normal bundle to the critical locus.
For the Heisenberg algebra this determinant equals $1$ (no higher interactions),
so the localized integral reduces to the sum over strata of the restriction
$\beta|_{\mathrm{crit}} \cdot \omega|_{\mathrm{crit}}$.  This sum is precisely
the bar-cobar counit
$\varepsilon \colon \barB_k(\mathrm{Heis}) \otimes \Omega_k(\mathrm{Heis}^!)
\to \mathbf{k}$
evaluated at $(\beta, \omega)$.
\end{proof}

\begin{remark}[The Chriss--Ginzburg interpretation]
\label{rmk:chriss-ginzburg-upgrade}
Proposition~\ref{prop:bv-bar-cobar-betagamma} upgrades
Conjecture~\ref{conj:bv-integration-bar-cobar} from a heuristic to a geometric
statement: the BV path integral is the \emph{pushforward of the fundamental
class of\/ $\Steinberg_b$ along the projection to the base}.  This is exactly
the definition of the convolution product in Chriss--Ginzburg~\cite{CG97}.
The heuristic becomes a theorem once the geometric Steinberg
(Conjecture~\ref*{conj:geometric-steinberg}) is established: the convolution
algebra of $\Steinberg_b$ recovers the bar-cobar adjunction of Vol~I, and the
BV path integral is the geometric shadow of the counit map.
\end{remark}

\begin{corollary}[\ClaimStatusConjectured]
\label{cor:bv-euler-steinberg}
The full BV partition function
\[Z_{\mathrm{BV}}(T) \;=\; \sum_{k \geq 0}
\int_{\Steinberg_b^{(k)}} [\Steinberg_b^{(k)}]\]
is the Euler characteristic of the Steinberg object:
$Z_{\mathrm{BV}} = \chi(\Steinberg_b)$.
This gives a new proof of the relationship between BV quantization and Koszul
duality: both extract the same invariant --- the Euler class --- from the same
geometric object, the Steinberg correspondence.
\end{corollary}

\subsubsection{Observables and correlation functions}

\begin{theorem}[Observables = cohomology; \ClaimStatusProvedElsewhere{} \cite{CG17}]
\label{thm:observables-brst-cohomology}
Physical observables are:
\[\mathcal{O}_{\text{phys}} = H^0(Q_{\text{BRST}}) =
\ker(Q_{\text{BRST}})/\text{im}(Q_{\text{BRST}})\]

In the bar-cobar framework:
\[\mathcal{O}_{\text{phys}} \cong H^0(\bar{B}^{\text{ch}}(\mathcal{A}))\]
\end{theorem}

\begin{example}[Correlation functions]
An $n$-point correlation function:
\[\langle \mathcal{O}_1(z_1) \cdots \mathcal{O}_n(z_n) \rangle =
\int_{\overline{C}_n(X)} [\mathcal{O}_1 \otimes \cdots \otimes \mathcal{O}_n]
\cdot e^{S_{\text{int}}}\]

is computed by representing each $\mathcal{O}_i$ as a cocycle in $\bar{B}^{\text{ch}}(\mathcal{A})$ and applying the bar-cobar pairing with $e^{S_{\text{int}}}$ --- Gaiotto's prescription for holomorphic-topological observables.
\end{example}

\subsection{Dictionary}

\begin{remark}[Dictionary]
The BV-BRST formalism connects:

\begin{center}
\begin{tabular}{|l|l|l|}
\hline
\textbf{Physics} & \textbf{Math (Bar-Cobar)} & \textbf{Geometry} \\
\hline
BV complex & Bar complex & Log forms on $\overline{C}_n$ \\
BV bracket & Configuration space symplectic & Poisson structure \\
Master equation & $\dzero^2 = 0$ & Stokes / boundary residues \\
BV Laplacian & Cobar differential & Delta functions \\
Quantum master eq & Bar-cobar duality & Verdier/Koszul pairing \\
BRST operator & Residue extraction & OPE singularities \\
Gauge fixing & Lagrangian choice & Regularization \\
Observables & Cohomology & Physical states \\
Path integral & Bar-cobar pairing & Configuration integrals \\
4d $\to$ 2d reduction & Factorization & Dimensional analysis \\
W-algebra & Boundary chiral algebra & Higgs moduli \\
\hline
\end{tabular}
\end{center}
\end{remark}

\begin{remark}[Gaiotto's contribution]
Gaiotto's contribution was recognizing a physical pattern: holomorphic-topological theories naturally produce chiral algebras, boundary conditions for these theories are modules over the chiral algebra, the open-closed correspondence should be modeled by bar-cobar duality, and suitable 4d gauge theory reductions should produce W-algebras via this mechanism.  The geometric bar-cobar construction supplies the theorematic boundary algebraic package on which these physical comparison statements are built.
\end{remark}
